

\documentclass[11pt]{article}

\usepackage[T1]{fontenc}
\usepackage[utf8]{inputenc}
\usepackage[english]{babel}
\usepackage[margin=1in]{geometry}
\usepackage{lmodern}
\usepackage{microtype}
\usepackage{mathtools}
\usepackage{amsmath,amssymb,amsthm}
\usepackage{array}
\usepackage{booktabs}
\usepackage{tabularx}
\usepackage{graphicx}
\usepackage{pdflscape}
\usepackage{caption}
\usepackage{enumitem}
\usepackage{fancyhdr}
\usepackage[numbers]{natbib}
\usepackage{doi}
\usepackage{xurl}
\usepackage{hyperref}
\usepackage{cleveref}
\usepackage{orcidlink}
\usepackage{titling}
\usepackage{tikz}
\usetikzlibrary{arrows.meta,positioning,calc,decorations.pathreplacing,shapes.geometric,fit,backgrounds}
\usepackage{pgfplots}
\pgfplotsset{compat=1.18}



\newcommand{\HF}{\mathrm{HF}}                 \newcommand{\rank}{\operatorname{rank}}
\newcommand{\TC}{\operatorname{TC}}
\newcommand{\dom}{\operatorname{dom}}
\newcommand{\Fix}{\operatorname{Fix}}
\newcommand{\im}{\operatorname{im}}
\newcommand{\Ord}{\mathrm{Ord}}
\newcommand{\powset}{\mathcal{P}}
\newcommand{\clF}{\operatorname{cl}_F}
\newcommand{\dec}{\operatorname{dec}}
\newcommand{\Ch}{\operatorname{Ch}}
\newcommand{\Sig}{\Sigma}                      \newcommand{\restr}{\!\upharpoonright\!}
\newcommand{\Cgen}{\mathbb{P}}                 \newcommand{\Fn}{\operatorname{Fn}}
\newcommand{\Vfin}{V_{\omega}}

\newcommand{\gradetag}[1]{\textnormal{\textsc{[#1]}}}
\newcommand{\inhouse}{\gradetag{in-house}}
\newcommand{\condg}{\gradetag{conditional}}
\newcommand{\modelg}{\gradetag{model}}
\newcommand{\guardg}{\gradetag{guard}}
\newcommand{\metag}{\gradetag{meta}}
 
\newtheorem{theorem}{Theorem}[section]
\newtheorem{lemma}[theorem]{Lemma}
\newtheorem{proposition}[theorem]{Proposition}
\newtheorem{corollary}[theorem]{Corollary}
\theoremstyle{definition}
\newtheorem{definition}[theorem]{Definition}
\newtheorem{remark}[theorem]{Remark}
\newtheorem{convention}[theorem]{Convention}
\urlstyle{same}
\captionsetup{
  font=small,
  labelfont=bf,
  labelsep=period
}
\captionsetup[table]{skip=6pt,position=top}
\crefname{theorem}{theorem}{theorems}
\Crefname{theorem}{Theorem}{Theorems}
\crefname{lemma}{lemma}{lemmas}
\Crefname{lemma}{Lemma}{Lemmas}
\crefname{proposition}{proposition}{propositions}
\Crefname{proposition}{Proposition}{Propositions}
\crefname{corollary}{corollary}{corollaries}
\Crefname{corollary}{Corollary}{Corollaries}
\crefname{definition}{definition}{definitions}
\Crefname{definition}{Definition}{Definitions}
\crefname{remark}{remark}{remarks}
\Crefname{remark}{Remark}{Remarks}
\setlist[enumerate]{leftmargin=*, itemsep=2pt, topsep=4pt}
\setlength{\droptitle}{-3.2em}
\pretitle{\begin{center}\LARGE\bfseries}
\posttitle{\par\end{center}\vspace{0.5em}}
\preauthor{\begin{center}\normalsize}
\postauthor{\par\end{center}\vspace{-0.4em}}
\predate{\begin{center}\small}
\postdate{\par\end{center}\vspace{0.6em}}
\setlength{\emergencystretch}{2em}
\clubpenalty=10000
\widowpenalty=10000
\displaywidowpenalty=10000
\interfootnotelinepenalty=10000
\allowdisplaybreaks[2]
\fancypagestyle{firstpage}{\fancyhf{}
  \fancyfoot[C]{\scriptsize Automorph Inc., Wilmington, DE, USA \quad \texttt{ioannis@automorph.io} \quad \textcopyright\ 2026 \quad CC-BY 4.0 \quad DOI: \doi{10.5281/zenodo.23121590}}
  \renewcommand{\headrulewidth}{0pt}
  \renewcommand{\footrulewidth}{0pt}
}

\title{To Set a Stone with Six Birds:\\
The Free and Purchased Parts of ZFC}
\author{Ioannis Tsiokos\,\orcidlink{0009-0009-7659-5964}}
\date{v1: 17 June 2026\quad$\cdot$\quad v2: 3 October 2026}
\hypersetup{
  pdftitle={To Set a Stone with Six Birds: The Free and Purchased Parts of ZFC},
  pdfauthor={Ioannis Tsiokos},
  hidelinks
}

\begin{document}

\maketitle
\thispagestyle{firstpage}

\begin{abstract}
The axioms of $\mathsf{ZFC}$ are usually presented as a flat list. We show that the list splits
into two parts: axioms that a small finite construction supplies for free, and axioms that
must be bought with explicit commitments. A set is treated as a \emph{package}: a finite
membership graph is observed only through its membership profile, and a finite Mostowski
collapse replaces the part of the graph reachable from its distinguished node by a canonical
representative. On finite graphs this packaging map is
idempotent with exactly the hereditarily finite sets $\HF$ as its objects; closing the empty
family under the empty set, pairing, union, subsets and finite images yields exactly $\HF$; and $(\HF,\in)$
satisfies every axiom of $\mathsf{ZFC}$ except Infinity. Foundation coincides, on finite graphs,
with the existence of a strictly decreasing rank, and over $\HF$ rank is first-order
definable from membership, so forgetting the stages loses nothing. No finitely generated rule started from hereditarily finite sets ever produces an
infinite set: Infinity is a change of rules, not a longer run of the old ones. The infinitary
axioms then enter through four named commitments---Infinity, Power set, Replacement and
Choice---on top of a stated structural premise; under them the membership structure satisfies
$\mathsf{ZF}$, or $\mathsf{ZFC}$, and standard countermodels show that each commitment is needed.

The second theme is forcing. A counting lemma from the Six Birds framework---a random predicate
on $N$ points is definable from a $K$-block partition with probability $2^{-(N-K)}$---counts
failures of the same disagreement conditions that Cohen forcing uses. Over $\omega$ these
conditions become dense, so a Cohen real over a countable transitive model $M$ is new and strictly
enlarges the Boolean algebra of $M$-coded predicates. Novelty does not suffice for genericity: we
give a real that is new and splits every infinite $M$-coded set without being generic, and we
state the exact test---a real is generic exactly when it avoids every $M$-coded closed nowhere
dense set.

Selected finite cores are verified in Lean without external libraries; an automatic audit of
all project declarations finds no axioms beyond propositional extensionality and quotient
soundness, so the finite results used to price Choice do not use Choice. We prove no consistency
or independence theorems; every infinitary ingredient is a named hypothesis or a cited
classical result.
\end{abstract}
 
\medskip
\noindent\textbf{Keywords.} set theory; ZFC; hereditarily finite sets; closure operators; Mostowski collapse; Cohen forcing; genericity; Boolean algebras of predicates; Six Birds Theory; emergence calculus; Lean.

\section{Introduction}\label{sec:intro}

\subsection{A list, and a question}

The foundations of mathematics are usually met as a list. Extensionality, empty set, pairing,
union, power set, separation, infinity, replacement, foundation, choice: ten statements about
membership \citep{Zermelo1908, Fraenkel1922} from which, we are told, the rest of mathematics
follows. The list is adopted because it works. It carries analysis, algebra and topology, and it
has not produced a contradiction.

This paper asks where the list comes from. Not whether it is true or consistent, but whether its
entries are ten independent stipulations, or whether some of them are consequences of something
simpler, with only a few genuinely \emph{chosen}. Our answer is that the list sorts into two
columns. Some axioms are \emph{free}: they hold automatically of anything produced by a modest
finite construction. Others are \emph{purchased}: they mark the exact points at which that
construction must be told to go on past where it would stop on its own. The free column is
proved; the purchased column is stated as a small set of named hypotheses, each with a theorem
or countermodel showing why it cannot be dropped.

\subsection{The idea in brief}

The organizing vocabulary comes from Six Birds Theory (\textsc{SBT})~\citep{SBTFoundationsI}, a
framework for describing how stable objects arise when a complicated substrate is observed
through a limited interface. We use only a small part of it, and we use it as an organizing
instrument, not as a source of new set theory. Three roles matter here.
\begin{itemize}[itemsep=2pt,topsep=3pt]
\item A \emph{lens} records what an observer can see. For sets, the substrate is the class of
finite membership graphs, and the lens records only the membership profile of the distinguished
node: which sets are its members, which sets are theirs, and so on. Node names and duplicated
structure are invisible.
\item A \emph{packaging map} replaces each graph by a canonical representative of what the lens
sees. For sets it collapses the part of the graph reachable from the distinguished node, as in
a finite Mostowski collapse, and re-presents it canonically; its objects are exactly the
hereditarily finite sets. Extensionality is the decision to identify graphs that the lens cannot tell apart.
\item An \emph{audit} is a certificate that strictly decreases along the structure. For sets it
is the rank, which decreases along membership. On finite graphs, Foundation says precisely that
such a certificate exists.
\end{itemize}
Two further patterns from the framework organize the rest of the paper. A fixed finite rule,
iterated, may keep producing new objects, but they stay of the same kind: everything it produces
is hereditarily finite, and it never leaves $\HF$. In this sense it \emph{saturates}. An object
of a new kind, such as an infinite set, comes only from \emph{changing the rule}, that is, from
adjoining something the old rule cannot produce. The free axioms are what the finite rules
already supply inside $\HF$; the purchased axioms are the named changes of rule. Figure~\ref{fig:overview} shows the resulting picture of the cumulative
hierarchy.

\begin{figure}[tbp]
\centering
\begin{tikzpicture}[
    font=\small,
    box/.style={draw, rounded corners=2pt, align=left, font=\footnotesize,
                inner sep=4pt, text width=5.4cm, execute at begin node=\raggedright},
    >={Stealth[length=2mm]}]
\coordinate (O) at (0,0);
  \coordinate (L) at (-2.4,6.0);
  \coordinate (R) at (2.4,6.0);
  \coordinate (Lw) at (-0.96,2.4);
  \coordinate (Rw) at (0.96,2.4);
  \fill[blue!10] (O) -- (Lw) -- (Rw) -- cycle;
  \fill[orange!12] (Lw) -- (L) -- (R) -- (Rw) -- cycle;
  \draw[thick] (L) -- (O) -- (R);
  \draw[densely dashed, thick] ($(Lw)+(-0.3,0)$) -- ($(Rw)+(0.3,0)$);
  \node[font=\footnotesize, anchor=west] at ($(Rw)+(0.32,0)$) {$\omega$};
  \foreach \t in {0.3,0.55,0.78} {
    \draw[gray!70] ($(O)!\t!(Lw)$) -- ($(O)!\t!(Rw)$);
  }
  \node[align=center, font=\footnotesize] at (0,1.25) {$\HF$};
  \node[align=center, font=\footnotesize] at (0,4.3) {stages $V_\alpha$\\$\alpha\ge\omega$};
\node[box, fill=blue!7, anchor=east] (F) at (-1.4,1.9) {\textbf{Free} (Sections~\ref{sec:free}--\ref{sec:shadow})\\[2pt]
    Extensionality: the lens quotient\\
    Empty, Pairing, Union: packaging\\
    Separation: gates (subsets)\\
    Replacement: finite images\\
    Power set, Choice: finite forms\\
    Foundation: the rank audit};
  \draw[->, gray] (F.east) -- (-0.8,1.9);
\node[box, fill=orange!12, anchor=west] (B) at (2.3,4.4) {\textbf{Purchased} (Section~\ref{sec:purchases})\\[2pt]
    \textbf{C-INF}: admit an inductive set\\
    \textbf{C-POW}: collect all subsets\\
    \textbf{C-REPL}: unbounded images\\
    \textbf{C-SEL}: sections of all surjections};
  \draw[->, gray] (B.west) -- (1.55,4.4);
\node[box, draw=none, anchor=west] (N) at (2.3,1.5) {No finitely generated rule started\\in $\HF$ crosses the dashed line\\(Theorem~\ref{thm:nonattain})};
  \draw[->, gray] (N.west) -- (1.15,2.25);
\end{tikzpicture}
\caption{The free/purchased picture of the cumulative hierarchy. Below $\omega$ (blue), a finite
packaging construction supplies $\HF$, in which every axiom except Infinity holds; no finitely
generated rule iterated inside $\HF$ reaches the dashed line. Above it (orange), four named
commitments, on top of a structural premise that applies the free clauses at every stage
(Theorem~\ref{thm:shadow}), yield $\mathsf{ZF}$ and, with \textbf{C-SEL}, $\mathsf{ZFC}$.}
\label{fig:overview}
\end{figure}
 
\subsection{Main results}

We work first over the hereditarily finite sets $\HF$, where everything is finite and can be
checked, and then pass to the infinite by explicit hypotheses.

\paragraph{The free part (\S\ref{sec:free}).} On finite well-founded membership graphs, the
packaging map is idempotent and its fixed points correspond bijectively to $\HF$
(Theorem~\ref{thm:p5}). Closing the empty family under the finite packaging operations---empty
set, pairing, union, subsets and finite images---gives exactly $\HF$ (Theorem~\ref{thm:hfclosure}),
and $(\HF,\in)$ satisfies every axiom of $\mathsf{ZFC}$ except Infinity, each witnessed by a
named clause, by finite selection, or by the rank audit (Theorem~\ref{thm:hfsat}). On finite graphs, Foundation is
equivalent to the existence of a strictly decreasing rank function (Theorem~\ref{thm:tfae}).

\paragraph{Infinity is a change of rules (\S\ref{sec:infinity}).} No finitely generated
completion rule, iterated from hereditarily finite sets, ever produces an inductive set
(Theorem~\ref{thm:nonattain}). An infinite set has to be admitted; it cannot be reached.

\paragraph{Forgetting the stages is lossless (\S\ref{sec:shadow}).} The usual axioms mention
only membership, while the construction is staged. Over $\HF$ rank is first-order definable from
membership, and every sentence about the staged structure translates into an equivalent sentence
about membership alone (Theorem~\ref{thm:translation}).

\paragraph{The purchased part (\S\ref{sec:purchases}).} Infinity, Power set, Replacement and
Choice enter as four named commitments. Infinity is provably not free
(Theorem~\ref{thm:nonattain}). Choice is, over $\mathsf{ZF}$, exactly the statement that every
surjection of sets has a section (Theorem~\ref{thm:choice}); the particular membership lens needs
no choice, because its sections are given explicitly. Under a stated structural premise and the
four commitments, the membership structure satisfies $\mathsf{ZF}$, and with Choice $\mathsf{ZFC}$
(Theorem~\ref{thm:shadow}). Standard countermodels show that none of the four can be omitted
(Proposition~\ref{prop:nonredundancy}). Table~\ref{tab:factorization} records the resulting
ledger.

\paragraph{Forcing (\S\ref{sec:forcing}).} The Six Birds framework contains a counting
lemma---informally called the ``finite forcing lemma''~\citep{SBTFoundationsI}---stating that a
random predicate on $N$ points is constant on every block of a $K$-block partition with
probability $2^{-(N-K)}$. We show that this lemma counts the failures of the same disagreement
conditions that Cohen forcing uses~\citep{Cohen1963, Cohen1966}, and we compute exactly which
finite conditions have no extension meeting them (Theorem~\ref{thm:densitydefect}). Over
$\omega$ the failures disappear. As a consequence, a Cohen-generic real over a countable
transitive model $M$ lies outside $M$ and splits every infinite set coded in $M$
(Theorem~\ref{thm:novelty}), and adjoining it strictly enlarges the Boolean algebra of
$M$-coded predicates (Theorem~\ref{thm:strict}). The implication runs one way only: a real can
be new and split every infinite coded set without being generic
(Proposition~\ref{prop:converse}). The exact condition is that the real avoid every $M$-coded
closed nowhere dense set (Theorem~\ref{thm:recognition}). Cohen reals form a comeagre set of
measure zero, so the finite probability $1-2^{-(N-K)}$ describes novelty, not genericity
(\S\ref{ssec:category}). The continuum hypothesis appears as a question the ground model does
not settle (\S\ref{sec:ch}).

\paragraph{Guards and verification (\S\S\ref{sec:atlas}--\ref{sec:mech}).} Eight countermodels
block the most tempting overstatements, such as ``extensionality is forced'' or ``Foundation is
free without staging.'' Selected finite cores are verified in Lean~4 without external libraries.
An automatic audit, run on every build, checks that each of the 371 project declarations depends
on no axiom other than propositional extensionality and quotient soundness; in particular, the
finite lift used to price Choice is verified without Choice. Reproducible numerical exhibits
accompany the counting results.

\subsection{Scope and grades}

We do \emph{not} derive $\mathsf{ZF}$ or $\mathsf{ZFC}$ from finite principles: the infinitary
content enters only through named hypotheses. We prove no consistency statement beyond the
classical results we cite, and no new independence result; the forcing section reorganizes
classical facts and adds finite and structural statements about them. We do not claim that
$\mathsf{ZFC}$ is the unique or most natural set theory. The interpretive reading in
\S\ref{sec:discussion} comes after the theorems and is never used to prove them.

Because the paper mixes statements of different strength, every numbered result carries a tag.
\inhouse{} marks a result proved here without set-theoretic hypotheses beyond the ambient
metatheory; some of these have Lean-verified cores. \condg{} marks a result proved under named
hypotheses or cited classical imports. \modelg{} marks a statement about specific models obtained
from cited results. \guardg{} marks a countermodel. \metag{} marks an interpretive reading that is
deliberately not a theorem. Section~\ref{sec:scope} collects the boundaries in one place.
 \section{The packaging instrument}\label{sec:instrument}

This section fixes the object the rest of the paper studies: set theory read as a Six Birds
\emph{theory package}~\citep{SBTFoundationsI}. We keep the framework's apparatus to the minimum
we use. Every presentation below is a finite structure, so the finite-carrier arguments of the
framework (its regime \textbf{A\_FIN}) apply to each presentation separately. The collection of
all presentations is infinite, and we never treat it as a finite carrier.

\subsection{Packages, lenses, audits}

\begin{convention}[Theory package]\label{conv:package}
A \emph{theory package} is a tuple $(Z,\,f,\,E,\,\mathcal A)$ consisting of a carrier $Z$; a
\emph{lens} $f\colon Z\to X$ recording what an observer sees of a carrier object; a
\emph{packaging} endomap $E\colon Z\to Z$, typically idempotent, whose fixed points are the
\emph{objects} the package recognizes; and an \emph{audit} $\mathcal A$, a certificate that
strictly decreases along the structure being audited. A predicate on $Z$ is \emph{visible} to
the lens if it is constant on every fiber of $f$.
\end{convention}

\subsection{The set-theoretic instrument}

\begin{definition}[Finite membership presentations]\label{def:presentation}
A \emph{finite pointed membership presentation} is a triple $P=(V,\to,v)$ where $V$ is a finite
set of nodes, ${\to}\subseteq V\times V$ is acyclic, and $v\in V$ is a distinguished
\emph{point}. We read $u\to w$ as ``$w$ is a member-candidate of $u$'' and write
$\Ch_P(u):=\{w: u\to w\}$. To keep the carrier a set, node sets are finite subsets of
$\mathbb N$. The carrier $Z$ is the set of all such presentations; it is countably infinite.
\end{definition}

In a finite graph, acyclicity is the same as well-foundedness: an infinite descending path in a
finite graph must repeat a node and therefore contains a cycle. This identification is the seed
of Foundation, and we return to it in \S\ref{ssec:foundation}.

\begin{definition}[The membership lens]\label{def:lens}
For $P=(V,\to,v)$ define the \emph{decoration} $\dec_P\colon V\to\HF$ by well-founded
recursion,
\[
  \dec_P(u)\;:=\;\{\dec_P(w): w\in\Ch_P(u)\},
\]
and set $f\colon Z\to\HF$, $f(P):=\dec_P(v)$. We call $f$ the \emph{membership-profile
lens}.\footnote{The term \emph{decoration} follows the graph-decoration viewpoint on membership
structures of \citet{Aczel1988}; on well-founded graphs it is the map collapsed by the Mostowski
lemma.}
\end{definition}

\begin{lemma}[Decoration is well posed; \inhouse]\label{lem:L1}
For every $P\in Z$ the decoration $\dec_P$ is uniquely defined, and $f(P)\in\HF$.
\end{lemma}

\begin{proof}
Finiteness and acyclicity give each node a finite height $h(u)$, the length of the longest
path starting at $u$, with $h(w)<h(u)$ whenever $u\to w$. Define $\dec_P$ by induction on
height: nodes of height $0$ are childless and decorate to $\varnothing$, and a node of height $n$
decorates to the finite set of the decorations of its children, all of lower height. The same
induction shows that each decoration is hereditarily finite and that the decoration is unique.
\end{proof}

\begin{definition}[Canonical presentation and packaging map]\label{def:E}
For $a\in\HF$ let $\TC(\{a\})$ be the least transitive set containing $a$ as an element, and let
$C(a):=(\TC(\{a\}),\in,a)$ be the \emph{canonical presentation} of $a$: its nodes are the
elements of $\TC(\{a\})$, named by their Ackermann codes in $\mathbb N$~\citep{Ackermann1937},
its edges are membership, and its point is $a$. The \emph{packaging map} is
\[
  E\colon Z\to Z,\qquad E(P):=C(f(P)).
\]
\end{definition}

In words, $E$ reads the membership profile of a presentation, forgets the presentation's
accidental shape, and presents the profile again in canonical form. Figure~\ref{fig:packaging}
shows a small example.

\begin{figure}[tbp]
\centering
\begin{tikzpicture}[
    font=\small,
    nd/.style={circle, draw, fill=white, minimum size=5.5mm, inner sep=0pt},
    pt/.style={circle, draw, thick, fill=blue!12, minimum size=5.5mm, inner sep=0pt},
    dl/.style={font=\scriptsize, text=blue!60!black},
    e/.style={-{Stealth[length=1.8mm]}, thick},
    >={Stealth[length=2mm]}]
\begin{scope}
    \node[pt] (r) at (0,2.4) {$r$};
    \node[nd] (u1) at (-1.3,1.2) {$u_1$};
    \node[nd] (u2) at (-0.35,1.2) {$u_2$};
    \node[nd] (w) at (1.1,1.2) {$w$};
    \node[nd] (u3) at (1.1,0) {$u_3$};
    \draw[e] (r) -- (u1); \draw[e] (r) -- (u2); \draw[e] (r) -- (w); \draw[e] (w) -- (u3);
    \node[dl, left=1pt of u1] {$\varnothing$};
    \node[dl, below=1pt of u2] {$\varnothing$};
    \node[dl, right=1pt of w] {$\{\varnothing\}$};
    \node[dl, right=1pt of u3] {$\varnothing$};
    \node[dl, right=2pt of r] {$\{\varnothing,\{\varnothing\}\}$};
    \node[font=\footnotesize] at (0,-0.8) {presentation $P$ (5 nodes)};
  \end{scope}
\draw[->, very thick, gray] (2.9,1.3) -- node[above, font=\footnotesize, text=black]
        {$E=C\circ f$} node[below, font=\scriptsize, text=black, align=center]
        {read the profile,\\re-present it} (5.0,1.3);
\begin{scope}[xshift=7.2cm]
    \node[pt] (two) at (0,2.4) {$2$};
    \node[nd] (one) at (0.9,1.2) {$1$};
    \node[nd] (zero) at (0,0) {$0$};
    \draw[e] (two) -- (one); \draw[e] (two) -- (zero); \draw[e] (one) -- (zero);
    \node[font=\footnotesize, align=center] at (0,-0.8)
         {canonical $C(2)=E(P)$ (3 nodes)};
    \draw[->, gray, thick] (-0.25,2.75) arc (200:-20:0.32);
    \node[font=\scriptsize, anchor=south] at (0.05,3.05) {$E(E(P))=E(P)$};
  \end{scope}
\end{tikzpicture}
\caption{The packaging map on a small example. The point $r$ of $P$ (shaded) decorates to
$2=\{\varnothing,\{\varnothing\}\}$; the blue labels are the decorations $\dec_P$. The lens $f$ keeps only
this profile, and $E$ re-presents it as the canonical graph $C(2)$ on the transitive closure
$\TC(\{2\})=\{0,1,2\}$. The three childless nodes $u_1,u_2,u_3$ merge into the single node $0$.
Applying $E$ again changes nothing (Theorem~\ref{thm:p5}).}
\label{fig:packaging}
\end{figure}
 
\begin{definition}[The rank audit]\label{def:rank}
The audit is $\rho(P):=\rank(f(P))$, and for a node $u$ of $P$, $\rho_P(u):=\rank(\dec_P(u))$,
where $\rank(a)=\sup\{\rank(b)+1: b\in a\}$.
\end{definition}

\subsection{Extensionality as a choice of lens}

The membership lens identifies two presentations exactly when they have the same profile. This
identification is not a theorem about presentations; it is the choice of instrument.

\begin{lemma}[Lens equivalence; \inhouse]\label{lem:L2}
Define $P\sim Q$ iff $f(P)=f(Q)$. Then $\sim$ is an equivalence relation; each class contains the
canonical presentation of its profile, since $f(C(a))=a$; and on rooted extensional
presentations\footnote{Rooted: every node is reachable from the point. Extensional: distinct
nodes have distinct sets of children.} $P\sim Q$ holds exactly when $P$ and $Q$ are isomorphic as
pointed graphs.
\end{lemma}

\begin{proof}
The first clause is immediate. For $f(C(a))=a$, show $\dec_{C(a)}(s)=s$ for every
$s\in\TC(\{a\})$ by induction on rank. For the third clause, let $P$ be rooted and extensional.
The decoration $\dec_P$ is injective: if $\dec_P(u)=\dec_P(u')$, induct on $h(u)+h(u')$; equal
decorations give a matching between the children of $u$ and of $u'$ with equal decorations, the
inductive hypothesis identifies matched children, and extensionality then identifies $u$ and
$u'$. Rootedness makes the image of $\dec_P$ exactly $\TC(\{f(P)\})$, so $\dec_P$ is an
isomorphism of pointed graphs from $P$ onto $C(f(P))$. Two such presentations with the same
profile are therefore isomorphic to the same canonical presentation, and the converse is clear.
\end{proof}

\begin{remark}[What the lens forgets]\label{rem:lens-forgets}
Lemma~\ref{lem:L2} says that extensional identification is a genuine quotient: node names, the
multiplicity of routes, and graph shape beyond the profile are invisible to this instrument. In
Figure~\ref{fig:packaging} the three childless nodes are merged. Extensionality is therefore not
forced by the presentations themselves; it is the decision to look only through $f$. A different
instrument may lawfully keep apart what $f$ merges (countermodel \textbf{CM-1},
\S\ref{sec:atlas}).
\end{remark}

\subsection{The audit decreases along membership}

\begin{lemma}[Audit monotonicity; \inhouse]\label{lem:L3}
If $u\to w$ in $P$, then $\rho_P(w)<\rho_P(u)$. Equivalently, for $a\in\HF$ and $b\in a$,
$\rank(b)<\rank(a)$.
\end{lemma}

\begin{proof}
$\dec_P(w)\in\dec_P(u)$ by definition, and $\rank(a)=\sup\{\rank(b)+1: b\in a\}$.
\end{proof}

\paragraph{Role assignment.} In the framework's numbering of its six roles, the packaging map
$E$ is the packaging role (P5), whose fixed points are analysed in Theorem~\ref{thm:p5}; the rank
$\rho$ is the audit role (P6), whose exactness is Theorem~\ref{thm:tfae}; taking subsets of a
set (gating) is the constraint role (P2); forming images along functions is the descent or
transport role (P1); and the staging of sets by rank is the staging role (P4). The sixth role,
the route-mismatch diagnostic (P3), is not used to prove anything here.
 \section{The free part: \texorpdfstring{$\HF$}{HF} and the auditable axioms}\label{sec:free}

We now run the instrument from nothing and read off what it supplies. Three results carry the
section: the packaging map is an idempotent whose objects are the hereditarily finite sets
(\S\ref{ssec:p5}); closing the empty family under the finite packaging operations gives exactly
$\HF$, which satisfies every axiom of $\mathsf{ZFC}$ except Infinity (\S\ref{ssec:hf}); and on
finite graphs Foundation is the same as the existence of a rank audit (\S\ref{ssec:foundation}).
Table~\ref{tab:freepart} summarizes the outcome.

\subsection{The set universe as a packaging quotient}\label{ssec:p5}

\begin{theorem}[Mostowski collapse is packaging; \inhouse]\label{thm:p5}
The packaging map $E\colon Z\to Z$ of Definition~\ref{def:E} is idempotent, and
\[
  \Fix(E)\;=\;\im(E)\;=\;\{C(a): a\in\HF\}.
\]
The restricted lens $f\restr\Fix(E)\colon \Fix(E)\to\HF$ is a bijection with inverse $a\mapsto
C(a)$. With $i\colon\Fix(E)\hookrightarrow Z$ the inclusion and $r\colon Z\to\Fix(E)$,
$r(P):=E(P)$,
\[
  r\circ i=\mathrm{id}_{\Fix(E)},\qquad i\circ r=E,
\]
so $\Fix(E)$ is a system of representatives for the lens quotient $Z/{\sim}$.
\end{theorem}

\begin{proof}
Since $f(C(a))=a$ (Lemma~\ref{lem:L2}), $E(E(P))=C(f(C(f(P))))=C(f(P))=E(P)$. Idempotence gives
$\im(E)\subseteq\Fix(E)$, the reverse inclusion is trivial, and $\im(E)=\{C(a):a\in\HF\}$ by
definition. The restricted lens is surjective because $f(C(a))=a$ and injective because $P=E(P)=
C(f(P))$ on $\Fix(E)$. The splitting equations are immediate, and $[P]\mapsto E(P)$ is a
bijection from $Z/{\sim}$ onto $\Fix(E)$ because $E(P)\sim P$. This is the framework's
``idempotents split'' lemma~\citep{SBTFoundationsI} for the membership instrument; the
underlying collapse of a well-founded extensional relation onto a transitive set is the Mostowski
lemma~\citep{Mostowski1949}.
\end{proof}

\begin{remark}[The construction already in use]\label{rem:mathlib}
Theorem~\ref{thm:p5} describes how the universe of sets is commonly built in proof assistants.
The \texttt{PSet}/\texttt{ZFSet} development of \citet{Mathlib2020} defines a pre-set as a
type-indexed family of pre-sets---a well-founded membership tree with arbitrary branching---and a
set as a class of pre-sets under extensional equivalence. That is the quotient of presentations
by the lens relation $\sim$, and Theorem~\ref{thm:p5} is its finite, well-founded instance. The
Lean development accompanying this paper builds the finite version from scratch, as a quotient of
finite membership trees, without choosing representatives (\S\ref{sec:mech}).
\end{remark}

\begin{remark}[Finite Mostowski collapse]\label{rem:finitemostowski}
On rooted extensional presentations $E$ is an isomorphism (Lemma~\ref{lem:L2}); in general it
collapses. In Figure~\ref{fig:packaging} a five-node presentation is packaged to the three-node
canonical presentation of $2$. Nodes that cannot be reached from the point do not affect its
decoration and are discarded by $E$; on the part reachable from the point, $E$ is the finite case
of the Mostowski collapse. The infinitary statement is classical and is neither reproved nor used
here.
\end{remark}

\subsection{Closing from nothing yields exactly \texorpdfstring{$\HF$}{HF}}\label{ssec:hf}

By Theorem~\ref{thm:p5} we may work directly with the packaged objects, the elements of $\HF$. We
write the basic operations on $\HF$ as packaging clauses and show that closing the empty family
under them produces all of $\HF$.

\begin{definition}[The finite packaging operator]\label{def:F}
For $A\subseteq\HF$ let
\[
  \begin{aligned}
  F(A):={}&\{\varnothing\}\cup A
  \;\cup\;\underbrace{\{\{a,b\}:a,b\in A\}}_{\text{pairing (P5)}}
  \;\cup\;\underbrace{\{\textstyle\bigcup a:a\in A\}}_{\text{union (P5)}}\\
  &\cup\;\underbrace{\{S:S\subseteq a,\ a\in A\}}_{\text{gating (P2)}}
  \;\cup\;\underbrace{\{h[a]:a\in A,\ h\colon a\to A\}}_{\text{transport (P1)}},
  \end{aligned}
\]
where $h[a]=\{h(t):t\in a\}$. Because every $a\in\HF$ is finite, every subset $S\subseteq a$ and
every function $h\colon a\to A$ is admitted; no definability restriction is needed at this level.
Definability restrictions reappear, and matter, in the forcing section (\S\ref{sec:forcing}).
\end{definition}

\begin{lemma}[$\clF$ is a closure operator; \inhouse]\label{lem:closure}
$F$ is extensive and monotone, and $\HF$ is $F$-closed. Hence for each $A\subseteq\HF$ the least
$F$-closed superset $\clF(A)=\bigcap\{B: A\subseteq B\subseteq\HF,\ F(B)\subseteq B\}$ exists,
and $\clF$ is a closure operator: extensive, monotone and idempotent
\citep{DaveyPriestley2002,Tarski1955}.
\end{lemma}

\begin{proof}
Extensivity is built in. For monotonicity, if $A\subseteq B$ then every output of a clause applied
to $A$ is an output of the same clause applied to $B$; in particular a map $h\colon a\to A$ is
also a map into $B$. $\HF$ is closed clause by clause, since each clause outputs a finite set of
hereditarily finite sets. The intersection of $F$-closed families containing $A$ is again
$F$-closed, by monotonicity, and the closure properties follow.
\end{proof}

\noindent Although $\HF$ is countable, the lattice $\powset(\HF)$ in which this intersection is
taken is uncountable. Forming the intersection uses the ambient set-theoretic metatheory (with no
choice principle), not a finite computation. The same is true of every statement below about
$\HF$ as a whole: each object and each operation is finite, but the global statements are
statements of the metatheory.

\begin{lemma}[Finite collection; \inhouse]\label{lem:collection}
If $A$ is $F$-closed and $a_1,\dots,a_n\in A$, then $\{a_1,\dots,a_n\}\in A$.
\end{lemma}

\begin{proof}
Induction on $n$, using $x\cup\{y\}=\bigcup\{x,\{y\}\}$: the cases $n\le 1$ use the empty-set
and pairing clauses, and the inductive step pairs $\{a_1,\dots,a_n\}$ with $\{a_{n+1}\}$ and takes
the union.
\end{proof}

\noindent Finite collection is thus \emph{derived} from pairing and union, not assumed.

\begin{theorem}[$\HF$ is the saturation point of finite packaging; \inhouse]\label{thm:hfclosure}
$\clF(\varnothing)=\HF$.
\end{theorem}

\begin{proof}
$\clF(\varnothing)\subseteq\HF$ because $\HF$ is $F$-closed. For the reverse inclusion, let
$C=\clF(\varnothing)$ and induct on rank: $\varnothing\in C$, and if every member of $a\in\HF$ lies
in $C$, then $a=\{t_1,\dots,t_n\}$ is a finite collection of elements of $C$, so $a\in C$ by
Lemma~\ref{lem:collection}.
\end{proof}

\begin{corollary}[Finite power sets are reachable; \inhouse]\label{cor:finpow}
For $a\in\HF$, $\powset(a)\in\clF(\varnothing)$: each $S\subseteq a$ enters by the gate clause,
and $\powset(a)$, a finite collection of such sets, enters by Lemma~\ref{lem:collection}.
\end{corollary}

\noindent Corollary~\ref{cor:finpow} marks the boundary on which the paper turns. \emph{Finite}
power sets are free, a consequence of gating and finite collection. Collecting \emph{all}
subsets of an infinite set into one object is a different act, and it must be purchased
(\S\ref{sec:purchases}).

\begin{theorem}[$(\HF,\in)$ satisfies $\mathsf{ZFC}$ without Infinity; \inhouse]\label{thm:hfsat}
The structure $(\HF,\in)$ satisfies Extensionality, Empty set, Pairing, Union, Power set, every
instance of the Separation and Replacement schemas, Foundation and Choice, and it satisfies the
negation of Infinity.
\end{theorem}

\begin{proof}
$\HF$ is transitive, so for $a,b\in\HF$ the statements ``$a\in b$'' and ``$a\subseteq b$'' mean
the same inside $(\HF,\in)$ as outside, and every subset of an element of $\HF$ is again in
$\HF$. Extensionality follows from transitivity. Empty set, Pairing, Union and Power set are
supplied by the corresponding clauses of $F$ and by Corollary~\ref{cor:finpow}.
\emph{Separation:} for $a\in\HF$ and a formula $\varphi$ with parameters in $\HF$, the set
$\{x\in a:(\HF,\in)\models\varphi(x)\}$ is a subset of the finite set $a$, hence in $\HF$ by the
gate clause. Only truth in $(\HF,\in)$ is used; no absoluteness of unbounded formulas is needed.
\emph{Replacement:} if $\varphi$ defines in $(\HF,\in)$ a function on $a\in\HF$, its unique values
form a finite function $h\colon a\to\HF$, and $h[a]\in\HF$ by the transport clause.
\emph{Foundation:} a nonempty $a\in\HF$ is finite, so it has an element of minimal rank, which is
$\in$-minimal by Lemma~\ref{lem:L3}. \emph{Choice:} a finite family of nonempty sets in $\HF$ has
a choice function in $\HF$; picking the member with the least Ackermann code is one explicit
choice, and the resulting finite set of ordered pairs is in $\HF$. \emph{Infinity fails:} let
$a\in\HF$ contain $\varnothing$, and let $x\in a$ have maximal rank. Then $x\cup\{x\}$ has rank
$\rank(x)+1$, so it is not in $a$, and $a$ is not inductive.
\end{proof}

\begin{remark}[Arithmetic]\label{rem:ackermann}
$(\HF,\in)$ is the classical model $\Vfin$ of $\mathsf{ZF}$ with Infinity replaced by its
negation \citep{Kunen2011}. Ackermann's coding \citep{Ackermann1937} gives mutually
inverse interpretations between $(\HF,\in)$ and the standard model of arithmetic. At the level of
theories, bi-interpretability with Peano arithmetic requires finite set theory to include, besides
the negation of Infinity, the statement that every set has a transitive closure
\citep{KayeWong2007}; $\HF$ satisfies that statement, so the remark applies to the model
studied here.
\end{remark}

\subsection{Foundation is rank-audit exactness}\label{ssec:foundation}

The axiom of Foundation goes back to the well-foundedness analyses of \citet{Mirimanoff1917} and
\citet{VonNeumann1928}. On finite graphs it is equivalent to the existence of an audit.

\begin{theorem}[Finite rank-audit equivalence; \inhouse]\label{thm:tfae}
For a finite directed graph $G=(V,\to)$ the following are equivalent:
\begin{enumerate}\itemsep0pt
  \item $G$ is acyclic (it has no directed cycle, and in particular no self-loop);
  \item there is an audit $\rho\colon V\to\mathbb N$ that strictly decreases along edges:
  $u\to w$ implies $\rho(w)<\rho(u)$;
  \item every nonempty $S\subseteq V$ has a $\to$-minimal element.
\end{enumerate}
\end{theorem}

\begin{proof}
$(1)\Rightarrow(2)$: in an acyclic graph every path has at most $|V|-1$ edges, so the height
of Lemma~\ref{lem:L1} (the length of the longest path from a node) is defined and strictly
decreases along edges. $(2)\Rightarrow(1)$: a cycle $u_0\to\cdots\to u_{k-1}\to u_0$ would give
$\rho(u_0)<\rho(u_0)$; a self-loop is the case $k=1$. $(2)\Rightarrow(3)$: an element of $S$
with least $\rho$-value is $\to$-minimal. $(3)\Rightarrow(1)$: the node set of a cycle is a
nonempty set without a $\to$-minimal element.
\end{proof}

\noindent Natural-number ranks always suffice for finite acyclic graphs. Infinite well-founded
relations may require ordinal-valued ranks.

\begin{corollary}[\inhouse]\label{cor:foundation}
On a canonical presentation $C(a)$ the rank is a strictly decreasing audit (Lemma~\ref{lem:L3}),
so by Theorem~\ref{thm:tfae} every nonempty $b\subseteq\TC(\{a\})$ has an $\in$-minimal element.
Conversely, any finite presentation that admits a strictly decreasing audit is acyclic.
\end{corollary}

\begin{remark}[Foundation as a certificate]\label{rem:foundation}
Theorem~\ref{thm:tfae} reads the prohibition of membership cycles as a certificate that a
staging discipline is in force. Membership graphs with cycles---the self-loop $x=\{x\}$, the
two-cycle $x=\{y\}$, $y=\{x\}$---are perfectly coherent finite objects. They are excluded by
\emph{this} instrument because they carry no audit; anti-foundation \citep{Aczel1988} admits them.
This is countermodel \textbf{CM-2} (\S\ref{sec:atlas}).
\end{remark}

\subsection{The free-part ledger}

\begin{table}[tbp]\centering
\small
\begin{tabular}{@{}lll@{}}
\toprule
Axiom in $(\HF,\in)$ & Supplier & Role \\
\midrule
Extensionality & lens quotient (Theorem~\ref{thm:p5}) & packaging boundary \\
Empty set & seed clause $\{\varnothing\}$ & P5 \\
Pairing & pairing clause & P5 \\
Union & union clause & P5 \\
Separation & gate clause & P2 \\
Power set (of finite sets) & gates $+$ finite collection (Corollary~\ref{cor:finpow}) & P2 $+$ P5 \\
Replacement (on finite domains) & transport clause & P1 \\
Choice (for finite families) & least-code selection & --- \\
Foundation & rank audit (Theorem~\ref{thm:tfae}) & P6 \\
$\neg\,$Infinity & every element is finite (Theorem~\ref{thm:hfsat}) & boundary \\
\bottomrule
\end{tabular}
\caption{The free part. Each axiom holds in $(\HF,\in)$ and is witnessed by a named clause of the
finite packaging operator $F$, by finite selection, or by the rank audit; the closure of the
empty family under these clauses is exactly $\HF$ (Theorem~\ref{thm:hfclosure}). All objects and
operations are finite; global statements about $\HF$ are made in the ambient metatheory.}
\label{tab:freepart}
\end{table}

Every entry of Table~\ref{tab:freepart} is proved here, and several finite cores behind the
entries are verified in Lean; the schema-level satisfaction statements remain on paper
(\S\ref{sec:mech}). The next two sections show that the boundary of this column is sharp:
Infinity cannot be reached from inside it (\S\ref{sec:infinity}), and forgetting the staging
loses nothing (\S\ref{sec:shadow}).
 \section{Infinity as a change of rules}\label{sec:infinity}

Theorem~\ref{thm:hfsat} left one axiom out of the free column: Infinity. This section shows that
the omission is forced. No fixed finite packaging rule, however often it is iterated from
hereditarily finite sets, produces an infinite set. Infinity is therefore the first place where
the construction must be told to continue: it is a change of rule, not a longer use of the old
one.

\subsection{Finite packaging rules}

The result is informative only if the class of excluded rules is broad enough to contain the
operations of \S\ref{sec:free}. Those operations draw new members not just from the current
family $A$ but also from inside its elements (the union and gate clauses). The right class is
therefore the following.

\begin{definition}[Finitely generated completion rule]\label{def:fg}
An operator $G$ on families of sets is a \emph{completion rule} if $A\subseteq G(A)$ for every
$A$. It is \emph{finitely generated} if every $x\in G(A)\setminus A$ is a \emph{finite} set with
\[
  x\;\subseteq\;A\cup\bigcup_{a\in A}\TC(a).
\]
\end{definition}

\begin{lemma}[The packaging operator is finitely generated; \inhouse]\label{lem:Ffg}
Extend $F$ (Definition~\ref{def:F}) to arbitrary families by requiring the union, gate and
transport outputs to be finite. The extended operator is finitely generated and agrees with $F$
on every $A\subseteq\HF$, where the finiteness requirement holds automatically.
\end{lemma}

The finiteness requirement is essential. Without it, the extension could form the union of an
infinite set, or one of its infinite subsets, and nothing would confine it to $\HF$.

\begin{lemma}[$\HF$-stability; \inhouse]\label{lem:hfstable}
If $G$ is finitely generated and $A\subseteq\HF$, then $G(A)\subseteq\HF$. Hence every finite
iterate $G^n(A_0)$ of a seed $A_0\subseteq\HF$ stays inside $\HF$.
\end{lemma}

\begin{proof}
A new element $x$ is a finite set whose members lie in $A$ or in the transitive closure of a
member of $A$. The transitive closure of a set in $\HF$ consists of sets in $\HF$, so every member
of $x$ is in $\HF$, and so is $x$. Induct on $n$.
\end{proof}

\subsection{Non-attainment}

\begin{theorem}[No finitely generated iteration reaches an inductive set;
\inhouse]\label{thm:nonattain}
Let $G$ be a finitely generated completion rule and $A_0\subseteq\HF$. Put $A_{n+1}=G(A_n)$ and
$A_\omega=\bigcup_n A_n$. Then:
\begin{enumerate}\itemsep0pt
  \item no $A_n$ contains an inductive set as an element;
  \item $A_\omega\subseteq\HF$, so $A_\omega$ does not contain $\omega$ or any other inductive
  set either; admitting $\omega$ requires a step outside the finitely generated class;
  \item if $G$ is moreover idempotent, the iteration stabilizes after one step:
  $A_n=A_1$ for all $n\ge1$.
\end{enumerate}
\end{theorem}

\begin{proof}
By Lemma~\ref{lem:hfstable} every $A_n$ is a subfamily of $\HF$, and no element of $\HF$ is
inductive (Theorem~\ref{thm:hfsat}). This gives (1), and (2) follows because a union of
subfamilies of $\HF$ is a subfamily of $\HF$; a step that added $\omega$ would add an infinite
set, which Definition~\ref{def:fg} forbids. (3) is immediate from $G(G(A_0))=G(A_0)$.
\end{proof}

Theorem~\ref{thm:nonattain} is a statement about objects produced by rules. It does not say that
finite axiom systems, or formal theories in general, cannot prove that an infinite set exists;
that is a different question, about syntax rather than construction.

\subsection{Infinity is a package change}

\begin{corollary}[The first purchase; \inhouse]\label{cor:packagechange}
Passing from $\HF$ to a family that contains an inductive set cannot be achieved by iterating any
finitely generated completion rule from any seed inside $\HF$. In the framework's terms, admitting
$\omega$ is a \emph{package change}: the new object is not produced by running the fixed package
longer.
\end{corollary}

We call the corresponding commitment \textbf{C-INF}: the admission of a completed infinite
collection, in the package-change sense of \citet{SBTIncompleteness}. We do not assert, derive
or evaluate \textbf{C-INF} here. We have shown only that it cannot be avoided by iteration.

\begin{remark}[Saturation and incompleteness; \metag]\label{rem:saturation}
Corollary~\ref{cor:packagechange} is one instance of a general pattern: a fixed completion rule
that adds only finite objects may keep adding them forever, but by iterating itself it cannot
produce an object of a new kind, here an infinite set. There
is a deliberate resonance with incompleteness. G\"odel's theorems \citep{Godel1931} say that a
consistent, effectively axiomatized theory containing enough arithmetic does not prove every
arithmetical truth, and (under the usual provability conditions) does not prove its own
consistency. Theorem~\ref{thm:nonattain} says that a fixed finite packaging rule does not reach
the infinite. Both describe a closed apparatus, run to saturation, with something deliberately out
of reach. The other half of the picture---what adjoining new content does---is the subject of
\S\ref{sec:forcing}. The resemblance between the two halves is a reading, not a metatheorem.
\end{remark}
 \section{The rank-erased shadow}\label{sec:shadow}

The construction so far carries two kinds of data: sets, and the stages at which they appear.
This is the \emph{iterative conception} of set in explicitly staged form
\citep{Boolos1971,Scott1974,Potter2004}. The axioms of $\mathsf{ZFC}$, by contrast, speak of a
single relation $\in$ and never mention stages. A natural worry is that the flat presentation
forgets the staging on which our reading of the axioms rests. This section shows that, over
$\HF$, nothing is forgotten: the stages are recoverable from membership alone.

\subsection{The staged structure and its shadow}

\begin{definition}[Staged structure and rank-erasure]\label{def:staged}
Let $L_2$ be the two-sorted language with sorts $\mathrm{SET}$ and $\mathrm{STAGE}$, a relation
$\in$ on sets, a relation $<$ on stages, and a function symbol $\mathrm{st}\colon\mathrm{SET}\to
\mathrm{STAGE}$. The canonical finite staged structure is
\[
  M_2:=(\HF;\ \Ord_{\mathrm{fin}};\ \in;\ <;\ \mathrm{st}=\rank),
\]
where $\Ord_{\mathrm{fin}}$ is the set of finite von Neumann ordinals, $<$ is membership between
ordinals, and $\mathrm{st}(x)=\rank(x)$. Its \emph{rank-erasure} is the one-sorted structure
$M_1:=(\HF,\in)$.
\end{definition}

\subsection{Rank is first-order recoverable}

\begin{lemma}[Finite ordinals are definable; \inhouse]\label{lem:orddef}
The formula $\mathrm{Ord}(u):\equiv\mathrm{Trans}(u)\wedge\mathrm{Lin}_\in(u)$, stating that $u$
is transitive and that any two distinct elements of $u$ are $\in$-comparable, defines over
$(\HF,\in)$ exactly the finite von Neumann ordinals.
\end{lemma}

\begin{proof}
Every finite ordinal is transitive and $\in$-linear. Conversely, induct on $|u|$. Because $\HF$ has
no membership cycles (Theorem~\ref{thm:tfae}), $\in$ is irreflexive on $u$, and it is transitive
on $u$: if $x\in y\in z$ in $u$, linearity leaves $x\in z$, $x=z$ or $z\in x$, and the last two
would create a cycle. So $\in$ is a strict linear order on the finite set $u$. If $u\ne\varnothing$,
let $m$ be its largest element. Every other element of $u$ lies in $m$, and $m\subseteq u$ by
transitivity of $u$, so $u=m\cup\{m\}$. Moreover $m$ is transitive (if $y\in x\in m$, then $y\in
u$ and $y\in m$ by transitivity of the order) and $\in$-linear, so by induction $m$ is a finite
ordinal and so is $u$.
\end{proof}

\begin{theorem}[Rank is definable; \inhouse]\label{thm:rankdef}
There is an $\{\in\}$-formula $R(x,r)$ such that, over $(\HF,\in)$, $R(x,r)$ holds iff
$r=\rank(x)$.
\end{theorem}

\begin{proof}
The witnesses of a rank computation are themselves hereditarily finite, so the recursion can be
quantified inside $\HF$. Write, as $\{\in\}$-formulas, Kuratowski pairs, functions and domains.
Call $(T,g)$ a \emph{rank attestation} for $x$ if $T$ is a transitive set with $x\in T$, $g$ is a
function with domain $T$, and
\[
  g(y)=\textstyle\bigcup_{z\in y}\bigl(g(z)\cup\{g(z)\}\bigr)\qquad\text{for all }y\in T.
\]
Let $R(x,r):\equiv\exists T\,\exists g\,\bigl((T,g)\text{ is a rank attestation for }x\wedge
g(x)=r\bigr)$. For existence take $T=\TC(\{x\})$ and $g=\rank\restr T$, both in $\HF$; an empty $y$
gets rank $0$. For uniqueness, any attestation satisfies $g(y)=\rank(y)$ for all $y\in T$, by
induction on rank. Hence $R(x,r)$ holds exactly when $r=\rank(x)$.
\end{proof}

\subsection{Lossless erasure}

\begin{theorem}[Rank-erasure translation; \inhouse]\label{thm:translation}
There is an effective translation $\Phi\mapsto\Phi^\flat$ from $L_2$-sentences to
$\{\in\}$-sentences such that
\[
  M_2\models\Phi\quad\Longleftrightarrow\quad M_1\models\Phi^\flat.
\]
Consequently the shadow $(\HF,\in)$ determines the stage sort, the stage order and the rank map.
\end{theorem}

\begin{proof}
First unnest locally at atoms: an atomic formula $\psi$ containing a term $\mathrm{st}(x)$ is
replaced by $\exists\alpha\,(\mathrm{st}(x)=\alpha\wedge\psi[\alpha/\mathrm{st}(x)])$ with $\alpha$
fresh. Because $\mathrm{st}$ is total and single-valued, this is equivalent to the atom, and since
the replacement happens inside each atom it commutes with negation and the other connectives. On
unnested formulas translate by recursion: $\mathrm{SET}$ variables and $\in$ are kept,
$\mathrm{STAGE}$ quantifiers are relativized to $\mathrm{Ord}$ (Lemma~\ref{lem:orddef}),
$(\alpha<\beta)^\flat:\equiv\alpha\in\beta$, $(\alpha=\beta)^\flat:\equiv\alpha=\beta$, and
$(\mathrm{st}(x)=\alpha)^\flat:\equiv R(x,\alpha)$ (Theorem~\ref{thm:rankdef}). An induction on
formulas, carrying variable assignments, shows that each formula and its translation agree; the
atomic cases are the two lemmas, and $\mathrm{STAGE}$ quantifiers range over the finite ordinals on
both sides.
\end{proof}

\begin{remark}[Scope of losslessness]\label{rem:translation}
Theorem~\ref{thm:translation} licenses the one-sorted bookkeeping of the rest of the paper: when
we say that the membership structure ``satisfies $\mathsf{ZF}$,'' no staging content has been
silently dropped. The infinitary analogue, the definability of the von Neumann rank in
$\mathsf{ZF}$, is standard \citep[Ch.~6]{Jech2003}, and deriving the axioms from the stage picture
is the classical complement of this recovery \citep{Scott1974,Boolos1971}.

Losslessness is a property of the canonical rank labeling, not of arbitrary stage labelings.
There are only countably many formulas, even with finitely many parameters from $\HF$, but
uncountably many maps $\HF\to\Ord_{\mathrm{fin}}$, so most labelings are not definable and would
be lost under erasure. The set-theoretic staging is exactly one that is not.
\end{remark}
 \section{The four commitments and the conditional shadow}\label{sec:purchases}

The free part stops at $\HF$. To obtain the usual universe of sets, and with it the full list of
$\mathsf{ZFC}$ axioms, the construction must be told to continue in four places. This section
states those four places as an explicit \emph{assumption box}, explains what each one costs, and
proves that under the box the membership structure satisfies $\mathsf{ZF}$, or $\mathsf{ZFC}$.
Everything infinitary enters through stated hypotheses; nothing in this section is presented as
derived from the finite construction. This follows the practice, used elsewhere in the series, of
stating a source of new content as a named hypothesis rather than building it into a definition
\citep{SBTRecognition,SBTIncompleteness}.

\subsection{The setting and the assumption box}\label{ssec:box}

We work in an ambient metatheory satisfying $\mathsf{ZF}$ (with Choice only where stated), and
consider a transitive class $U$, the \emph{admitted universe}. Statements ``in $U$'' are
relativized to $U$; for a proper class $U$ this is done separately for each fixed formula, without
a truth predicate. The \emph{structural premise} asks that the free clauses of $F$ be applied at
every stage:

\begin{description}[leftmargin=1.2em]
\item[\textbf{(S)} Structural premise.] $U$ is transitive, $\varnothing\in U$, and for all
$x,y\in U$ the sets $\{x,y\}$ and $\bigcup x$ belong to $U$. Moreover, for every formula $\varphi$
with parameters in $U$ and every $x\in U$, the set $\{t\in x: (U,\in)\models\varphi(t)\}$ belongs
to $U$.
\end{description}

The last clause is the Separation schema read as a gate at every stage. It is the infinitary form
of the gate clause, but it is not a consequence of it: finite gates do not prove infinitary
Separation, so (S) is assumed, not derived. The four commitments are:

\begin{description}[leftmargin=1.2em]
\item[\textbf{C-INF} (inductive-stage completion).] Some $I\in U$ is inductive: $\varnothing\in
  I$, and $x\cup\{x\}\in I$ whenever $x\in I$. \emph{Package reading:} the limit of the successor
  process is admitted as a completed object. \emph{Classical name:} Infinity.
\item[\textbf{C-POW} (full subpackage completion).] For every $x\in U$ there is $p\in U$ whose
  elements are exactly the $y\in U$ with $y\subseteq x$. \emph{Package reading:} all gates over
  $x$ are collected into one object at the next stage. \emph{Classical name:} Power set.
\item[\textbf{C-REPL} (unbounded transport).] For every formula $\varphi$ with parameters in $U$
  and every $x\in U$: if each $t\in x$ has exactly one $y\in U$ with $(U,\in)\models\varphi(t,y)$,
  then the set of these $y$ belongs to $U$. \emph{Package reading:} transport along definable maps
  is allowed at every scale. \emph{Classical name:} Replacement.
\item[\textbf{C-SEL} (prototype selection).] Every surjection in $U$ between elements of $U$ has
  a section in $U$; equivalently, every partition of a set in $U$ into nonempty fibers has a
  family of fiber representatives (\emph{prototypes}) in $U$. \emph{Classical name:} Choice.
\end{description}

Two readings deserve care. \textbf{C-POW} asks for the subsets of $x$ that lie in $U$; asking for
every ambient subset is a stronger, separate demand. And \textbf{C-INF} asks for some inductive
set; the least one, $\omega$, is then obtained inside $U$ by Separation.

\subsection{What each commitment costs}

\paragraph{C-INF is provably not free.} By Theorem~\ref{thm:nonattain}, no finitely generated
rule reaches an inductive set. Its non-freeness is a theorem.

\paragraph{C-POW exceeds its finite form.} Corollary~\ref{cor:finpow} gives power sets of finite
sets for free. \textbf{C-POW} adds the collection of all subsets of an infinite set into one
object, which no finite clause does.

\paragraph{C-REPL exceeds its finite form.} The transport clause of $F$ gives Replacement on
finite domains (Theorem~\ref{thm:hfsat}). \textbf{C-REPL} allows images of infinite sets under
definable maps, with no bound on the ranks of the values.

\paragraph{C-SEL is the price of sections for all surjections.} A lens is a surjection onto what
it sees, and a choice of prototype in each fiber is a section of it. For the membership lens of
\S\ref{sec:instrument} no choice is needed: $a\mapsto C(a)$ is an explicitly defined section.
More generally, a surjection whose source is well-ordered has a section in $\mathsf{ZF}$, namely
the least preimage. What \textbf{C-SEL} buys is a section for \emph{every} surjection of sets,
with no such structure given. The resulting prototypes are ``canonical'' only in the sense that
they are supplied; Choice provides no natural or symmetry-invariant selection.

\begin{theorem}[Choice is the price of sections; finite part \inhouse, ZF equivalence standard]\label{thm:choice}
\emph{(Finite part.)} Every surjection $f\colon Z\to X$ with $Z$ finite has a section, obtained
without any choice principle. \emph{(Equivalence.)} Over $\mathsf{ZF}$, the following are
equivalent: \textup{(1)} the Axiom of Choice; \textup{(2)} every surjection between sets has a
section; \textup{(3)} every partition of a set into nonempty fibers has a family of prototypes;
\textup{(4)} every product of a set-indexed family of nonempty sets is nonempty.
\end{theorem}

\begin{proof}
Finite part: a finite set $Z$ admits an enumeration $z_1,\dots,z_n$; asserting that one exists and
fixing it for this one set uses no choice principle. Let $s(x)$ be the member of $f^{-1}(x)$ with
least index. Equivalence: $(1)\Rightarrow(2)$ by choosing from the fibers; $(2)\Rightarrow(1)$
because, for nonempty sets $(A_i)_{i\in I}$, the projection $\{(i,a):a\in A_i\}\to I$ is a
surjection and a section is a choice function; $(2)\Leftrightarrow(3)$ because a section is the
same thing as a family of fiber representatives; $(1)\Leftrightarrow(4)$ is standard
\citep{Jech1973}.
\end{proof}

\begin{corollary}[Price of C-SEL; \condg]\label{cor:choice}
If $U$ satisfies $\mathsf{ZF}$, then \textbf{C-SEL} holds in $U$ if and only if $U$ satisfies the
Axiom of Choice. Choice is thus the license to select prototypes for arbitrary lenses: free for
finite lenses and for lenses that come with a well-ordering or an explicit section, and purchased
in general.
\end{corollary}

The equivalence concerns surjections between \emph{sets}. It says nothing about surjections whose
source or fibers are proper classes, and it is not used for them.

\subsection{The conditional shadow theorem}

\begin{theorem}[Conditional $\mathsf{ZF(C)}$ shadow; \condg]\label{thm:shadow}
Work in an ambient $\mathsf{ZF}$ and let $U$ be a transitive class satisfying \textup{(S)},
\textbf{C-INF}, \textbf{C-POW} and \textbf{C-REPL}. Then $(U,\in)$ satisfies every axiom of
$\mathsf{ZF}$, each schema instance by instance. If $U$ also satisfies \textbf{C-SEL}, then
$(U,\in)$ satisfies $\mathsf{ZFC}$.
\end{theorem}

\begin{proof}
\emph{Free suppliers.} Extensionality holds because $U$ is transitive. Empty set, Pairing, Union
and each instance of Separation are the clauses of (S). Foundation holds because a nonempty
$x\in U$ has an $\in$-minimal element in the ambient universe, and that element lies in $U$ by
transitivity; this uses the well-foundedness of the ambient membership relation, the infinitary
counterpart of the rank audit. \emph{Purchased suppliers.} Because $U$ is transitive and closed
under pairs and unions, ``$x\cup\{x\}$'' and ``inductive'' mean the same in $U$ as outside, so
\textbf{C-INF} gives Infinity. \textbf{C-POW} gives Power set, read in $U$, and \textbf{C-REPL}
gives each instance of Replacement. For Choice, let $x\in U$ be a set of nonempty sets. The set
$\{(i,a): i\in x,\ a\in i\}$ is in $U$, being a definable subset of $x\times\bigcup x$, which is
available from Power set and Separation. Its projection to $x$ is a surjection in $U$, and a
section supplied by \textbf{C-SEL} is a choice function.
\end{proof}

Inside such a $U$, the cumulative hierarchy $V_0=\varnothing$, $V_{\alpha+1}=\powset(V_\alpha)$,
$V_\lambda=\bigcup_{\beta<\lambda}V_\beta$ is defined by transfinite recursion \citep{Zermelo1930}
and every set has a definable rank \citep[Ch.~6]{Jech2003}. So $(U,\in)$ is the rank-erased
membership structure of its own staged hierarchy, the infinitary analogue of
Theorem~\ref{thm:translation}. These classical facts are imported, not reproved.

Theorem~\ref{thm:shadow} is a conditional statement about a given $U$; it does not produce one.
In particular, nothing here implies the existence of an inaccessible cardinal, of a transitive set
model of $\mathsf{ZF}$, or the consistency of $\mathsf{ZFC}$. For a truncated universe $V_\kappa$,
for instance, Replacement holds only if definable images of sets in $V_\kappa$ stay of rank below
$\kappa$, which an arbitrary cutoff does not guarantee \citep{MontagueVaught1959}.

\begin{table}[tbp]\centering
\small
\begin{tabular}{@{}llll@{}}
\toprule
Axiom & In $\HF$ (free) & At infinite stages & Why needed \\
\midrule
Extensionality & lens quotient (Thm.~\ref{thm:p5}) & transitivity of $U$ & --- \\
Empty, Pairing, Union & packaging clauses & premise (S) & --- \\
Separation & gate clause & premise (S) & --- \\
Foundation & rank audit (Thm.~\ref{thm:tfae}) & ambient well-foundedness & --- \\
\midrule
Infinity & fails (Thm.~\ref{thm:hfsat}) & \textbf{C-INF} & $\HF$; Thm.~\ref{thm:nonattain} \\
Power set & finite sets (Cor.~\ref{cor:finpow}) & \textbf{C-POW} & $H_{\omega_1}$ \\
Replacement & finite domains & \textbf{C-REPL} & $V_{\omega+\omega}$ \\
Choice & finite families & \textbf{C-SEL} & symmetric models of $\mathsf{ZF}$ \\
\bottomrule
\end{tabular}
\caption{The free/purchased ledger. The upper block holds in $\HF$ by finite clauses and holds at
infinite stages by the structural premise (S), which applies the same clauses at every stage, or
by ambient well-foundedness. The lower block is supplied at infinite stages only by the four named
commitments; the last column names the countermodel or theorem showing that the commitment cannot
be dropped (Proposition~\ref{prop:nonredundancy}).}
\label{tab:factorization}
\end{table}

\subsection{Non-redundancy}\label{ssec:nonredundancy}

The ledger is only meaningful if each purchase buys something the others do not.

\begin{proposition}[Each commitment is needed; mixed grades]\label{prop:nonredundancy}
\begin{enumerate}[label=\textup{(\roman*)}]
\item \emph{Dropping \textbf{C-INF}} \inhouse. $U=\HF$ satisfies (S), \textbf{C-POW},
\textbf{C-REPL} and \textbf{C-SEL}, and contains no inductive set (Theorem~\ref{thm:hfsat}).
\item \emph{Dropping \textbf{C-REPL}} \condg, in ambient $\mathsf{ZFC}$. $U=V_{\omega+\omega}$
satisfies (S), \textbf{C-INF}, \textbf{C-POW} and \textbf{C-SEL}, so it is a model of Zermelo set
theory with Choice. \textbf{C-REPL} fails: the map $n\mapsto\omega+n$ is definable in $U$ on the
domain $\omega\in U$, with each value in $U$, but its image has rank $\omega+\omega$ and is not an
element of $U$ \citep[Ch.~12]{Jech2003}.
\item \emph{Dropping \textbf{C-POW}} \condg, in ambient $\mathsf{ZFC}$. Let $U=H_{\omega_1}$, the
hereditarily countable sets. It satisfies (S), \textbf{C-INF}, \textbf{C-REPL} (indeed the stronger
Collection schema) and \textbf{C-SEL} \citep{GitmanHamkinsJohnstone2016}. \textbf{C-POW} fails for
$x=\omega$: every subset of $\omega$ is hereditarily countable and hence in $U$, so a power set of
$\omega$ in $U$ would contain all of the uncountably many reals, while every element of $U$ is
countable.
\item \emph{Dropping \textbf{C-SEL}} \condg, relative consistency. If $\mathsf{ZF}$ is consistent,
so is $\mathsf{ZF}$ together with the existence of a surjection that has no section; Cohen's
symmetric models give such models of $\mathsf{ZF}$ without atoms \citep{Cohen1966,Jech1973}. By
Corollary~\ref{cor:choice}, \textbf{C-SEL} fails in them.
\end{enumerate}
\end{proposition}

\begin{proof}[Proof notes]
(i) $\HF$ is transitive and closed under the clauses of (S) by Theorem~\ref{thm:hfsat}; it
satisfies Power set, Replacement and Choice internally by the same theorem, and these are exactly
\textbf{C-POW}, \textbf{C-REPL} and \textbf{C-SEL} read in $U=\HF$. (ii) Finite rank increments of
sets in $V_{\omega+\omega}$ stay below the limit $\omega+\omega$, which gives the closure
properties; a choice function of a set in $V_{\omega+\omega}$ has only finitely larger rank. (iii)
For Collection over a countable domain, pick one hereditarily countable witness per element using
ambient Choice; the countable union of their countable transitive closures is countable, by the
regularity of $\omega_1$. Each set in $U$ has a countable well-order in $U$, which gives internal
sections.
\end{proof}

The witness in (i) matters most for the paper's claim, because it is the boundary of the free
column itself: drop the first purchase and the remaining three are still satisfied, by $\HF$. The
other witnesses are classical; they are recalled here, not reproved.
 \section{Forcing as strict extension}\label{sec:forcing}

The Six Birds framework proves, in its finite setting, a counting lemma about how rarely a random
predicate is definable from a coarse partition. It was called, half in jest, the ``finite forcing
lemma''~\citep{SBTFoundationsI}, with an explicit warning that its kinship with Cohen's forcing
\citep{Cohen1963,Cohen1966} was only an analogy. This section makes the kinship precise and marks
where it stops.

The finite lemma counts the failures of the same \emph{disagreement} and \emph{splitting}
conditions that Cohen forcing uses (\S\ref{ssec:finitebridge}). Over $\omega$ those failures
disappear (\S\ref{ssec:vanish}), and as a consequence a Cohen-generic real over a countable
transitive model is new, splits every infinite set the model can code, and strictly enlarges the
algebra of predicates the model can express (\S\S\ref{ssec:genericity}--\ref{ssec:strict}). The
implication runs in one direction only. Novelty and splitting do not make a real generic; we give
an explicit counterexample and the exact test that does characterize genericity
(\S\ref{ssec:converse}). Finally, the finite probabilities measure novelty, not genericity:
Cohen reals form a set that is large in the sense of category but has measure zero
(\S\ref{ssec:category}).

\subsection{The finite Cohen bridge}\label{ssec:finitebridge}

Fix a finite index set $Z$ of size $N$ and a lens $f\colon Z\to X$ with $K$ nonempty blocks
$B_x=f^{-1}(x)$; the degenerate case $N=K=0$ is allowed. A \emph{predicate} is a total map
$h\colon Z\to\{0,1\}$. It is \emph{definable from $f$}, or visible to $f$, if it is constant on
every block. This is the paper's observational convention, not formula definability in a
specified language \citep[for which see][]{Libkin2004}. Equivalently, $h$ is definable from $f$
if the refined lens $z\mapsto(f(z),h(z))$ induces the same partition as $f$.

The \emph{finite Cohen poset} $\Cgen_Z$ consists of the partial functions $p\colon
Z\rightharpoonup\{0,1\}$, with $q\le p$ when $q$ extends $p$. A total predicate $h$ determines the
\emph{restriction filter} $G_h=\{p:p\subseteq h\}$. For a definable $g$, and for a block $B_x$ with
at least two elements, consider
\[
\begin{aligned}
  D_g&=\{p:\ p(z)\ne g(z)\text{ for some }z\in\dom p\},\\
  E_x&=\{p:\ p(z)\ne p(z')\text{ for some }z,z'\in B_x\cap\dom p\}.
\end{aligned}
\]

\begin{lemma}[Meeting equivalences; \inhouse]\label{lem:meeting}
For every total $h$:
\begin{enumerate}\itemsep0pt
\item $G_h$ meets $D_g$ iff $h\ne g$. Hence $h$ is not definable from $f$ iff $G_h$ meets $D_g$
for every definable $g$.
\item $G_h$ meets $E_x$ iff $h$ is not constant on $B_x$. Hence $h$ splits every block with at
least two elements iff $G_h$ meets every $E_x$.
\end{enumerate}
\end{lemma}

The two conditions in the lemma are different. A non-definable $h$ need only split \emph{some}
block; splitting \emph{every} non-singleton block is stronger.

\begin{theorem}[Density defects; \inhouse]\label{thm:densitydefect}
For a family $\mathcal D$ of conditions, let $F(\mathcal D)=\{p:\text{no }q\le p\text{ lies in
}\mathcal D\}$ be the set of conditions that have no extension in $\mathcal D$. Then
\[
  F(D_g)=\{g\},\qquad
  F(E_x)=\{p:\ B_x\subseteq\dom p\text{ and }p\restr B_x\text{ is constant}\}.
\]
So the only condition with no extension in $D_g$ is the total condition $g$ itself, and the
conditions with no extension in $E_x$ are exactly those that have already fixed the whole block to
a constant.
\end{theorem}

\begin{proof}
For $D_g$: a condition that is not a restriction of $g$ already disagrees with $g$ somewhere, so
it lies in $D_g$. A proper restriction of $g$ leaves some coordinate $z$ free, and setting it to
$1-g(z)$ gives an extension in $D_g$. The total condition $g$ has no extension but itself, which
is not in $D_g$. For $E_x$: if $p$ fixes $B_x$ to a constant, no extension splits the block.
Otherwise, either $p$ already splits $B_x$; or $p$ assigns at least one bit on $B_x$ and leaves
another coordinate of $B_x$ free, which can be set to the opposite value; or $p$ assigns no bit on
$B_x$, and two free coordinates of $B_x$ can be set to $0$ and $1$.
\end{proof}

\begin{remark}[Genericity in the finite poset is trivial]\label{rem:finitegeneric}
In $\Cgen_Z$ a total condition has no proper extension, so a genuinely dense family contains
every total condition, and every restriction filter $G_h$ meets every dense family. The families
$D_g$ are not dense: $g$ has no extension in $D_g$. The finite statements in this subsection are therefore about meeting
these particular \emph{defective} families, not about genericity in the finite poset.
\end{remark}

\begin{theorem}[The counting lemma; \inhouse]\label{thm:counting}
Let $h$ be a uniformly random total predicate. Then
\begin{enumerate}\itemsep0pt
\item $h$ is definable from $f$, equivalently $G_h$ misses some $D_g$, with probability exactly
$2^{-(N-K)}$;
\item $h$ splits every block with at least two elements, equivalently $G_h$ meets every $E_x$, with
probability exactly $\prod_{|B_x|\ge2}\bigl(1-2^{1-|B_x|}\bigr)$, and hence with probability at
least $1-\sum_{|B_x|\ge2}2^{1-|B_x|}$.
\end{enumerate}
\end{theorem}

\begin{proof}
(1) By Lemma~\ref{lem:meeting}, $G_h$ misses $D_g$ exactly when $h=g$. There are $2^K$ definable
predicates, one bit per block, among $2^N$ predicates in all, giving $2^K/2^N$. (2) The blocks are
disjoint, so the events ``$h$ is constant on $B_x$'' are independent, each with probability
$2\cdot2^{-|B_x|}$. The product and the union bound follow.
\end{proof}

When $K=N$ every block is a singleton: every predicate is definable, and splitting every
non-singleton block holds vacuously. Theorem~\ref{thm:counting}(1) \emph{is} the framework's
finite forcing lemma. Its rarity $2^{-(N-K)}$ is the probability that $h$ lands exactly on the
defect $F(D_g)=\{g\}$ of one of the disagreement families. The lemma is therefore the statistics of
the density families that forcing uses, at a finite size where those families are defective.

\subsection{Over \texorpdfstring{$\omega$}{omega} the defects vanish}\label{ssec:vanish}

Let $\Cgen=\Fn(\omega,2)$ be the poset of finite partial functions $p\colon\omega
\rightharpoonup\{0,1\}$, ordered by extension: the finite Cohen poset over the infinite index
$\omega$. Define $D_g$ for any total $g\colon\omega\to\{0,1\}$, and $E_A$ for any $A\subseteq\omega$,
exactly as above.

\begin{corollary}[Infinite-index contrast; \inhouse]\label{cor:infinitecontrast}
In $\Fn(\omega,2)$, $D_g$ is dense for every total $g$, and $E_A$ is dense for every
\emph{infinite} $A\subseteq\omega$. A finite $A$ with at least two elements keeps its defect.
\end{corollary}

\begin{proof}
A finite condition leaves infinitely many coordinates free. One of them can be set against $g$,
and an infinite $A$ contains two free coordinates, which can be set to $0$ and $1$. A condition
that fixes a finite $A$ to a constant has no extension splitting $A$.
\end{proof}

The finite defect $\{g\}$ disappears because no finite condition is total over $\omega$, and the
splitting defect disappears because no finite condition fixes an infinite block. What was rare in
the finite count becomes dense, and density is what genericity meets.

\subsection{Generic reals are new}\label{ssec:genericity}

We use two named hypotheses.
\begin{description}[leftmargin=1.2em]
\item[\textbf{J1}.] There is a countable transitive model $M$ of $\mathsf{ZFC}$. Its countability
is external: inside $M$, the set of reals is uncountable.
\item[\textbf{J2}.] The classical generic-extension machinery \citep{Cohen1966,Kunen2011}: if $G$
is $\Cgen$-generic over $M$, there is a transitive model $M[G]$ of $\mathsf{ZFC}$ with
$M\subseteq M[G]$, $G\in M[G]$, and the same ordinals as $M$.
\end{description}
The results of this subsection and the next use \textbf{J1} only; \textbf{J2} is used only to name
the model in which the new real lives.

\begin{definition}[Coded predicates]\label{def:groundlens}
Let $\mathcal B_M:=\powset(\omega)\cap M$, the subsets of $\omega$ that belong to $M$; we identify
each with its characteristic function, an \emph{$M$-coded predicate} on $\omega$. Since $M$ is a
transitive model, $\mathcal B_M$ contains $\varnothing$ and $\omega$ and is closed under finite
Boolean operations; it is a Boolean algebra of predicates on $\omega$.
\end{definition}

\begin{lemma}[Few coded predicates; \condg, under \textbf{J1}]\label{lem:countable}
$\mathcal B_M$ is countable, while $2^\omega$ is uncountable.
\end{lemma}

This is the infinitary counterpart of $2^K$ visible predicates among $2^N$. The analogy has a
limit, discussed in \S\ref{ssec:strict}: $\mathcal B_M$ is not the family of predicates visible to
any partition of $\omega$.

\begin{lemma}[Generics exist; \condg, under \textbf{J1}]\label{lem:RS}
For every $p\in\Cgen$ there is a filter $G\ni p$ that meets every dense subset of $\Cgen$ lying in
$M$.
\end{lemma}

\begin{proof}
Since $M$ is transitive and computes $\omega$ correctly, the poset $\Cgen$ as computed in $M$ is
the actual poset, and density of a set in $M$ means the same inside $M$ as outside. Because $M$ is
countable, enumerate its dense subsets of $\Cgen$ as $D_0,D_1,\dots$ and fix an enumeration of
$\Cgen$. Build $p=p_0\ge p_1\ge\cdots$, with $p_{n+1}$ the first extension of $p_n$ in $D_n$ in that
enumeration, so that no choice principle is used. The upward closure of the chain is a filter
containing $p$ and meeting every $D_n$. This is the Rasiowa--Sikorski lemma
\citep{RasiowaSikorski1950}.
\end{proof}

\begin{theorem}[Generic reals are new; \condg, under \textbf{J1}]\label{thm:novelty}
Let $G$ be $\Cgen$-generic over $M$ and $x_G:=\bigcup G$. Then $x_G\colon\omega\to\{0,1\}$ is
total, $G$ is exactly the restriction filter $\{p:p\subseteq x_G\}$, and:
\begin{enumerate}\itemsep0pt
\item $x_G\ne y$ for every $y\in\mathcal B_M$; in particular $x_G\notin M$;
\item $x_G$ splits every infinite $A\in\mathcal B_M$: there are $m,n\in A$ with $x_G(m)\ne x_G(n)$.
\end{enumerate}
\end{theorem}

\begin{proof}
Conditions in a filter are compatible, so $x_G$ is a function. The sets $T_n=\{p:n\in\dom p\}$ are
dense and lie in $M$, so $G$ meets each of them and $x_G$ is total. Every $p\in G$ is a
restriction of $x_G$; conversely, a finite restriction of $x_G$ is contained in a common extension
of finitely many members of $G$, which lies in $G$, so it lies in $G$ by upward closure. For (1),
let $y\in\mathcal B_M$. The family $D_y$ is dense (Corollary~\ref{cor:infinitecontrast}) and lies in
$M$, being defined there from $y$. So $G$ meets $D_y$, and $x_G$ disagrees with $y$ somewhere. Every
real in $M$ is such a $y$, so $x_G\notin M$. For (2), let $A\in M$ be infinite. Then $E_A$ is dense
(Corollary~\ref{cor:infinitecontrast}) and lies in $M$, so $G$ meets it.
\end{proof}

The generic real thus escapes the ground model through the same disagreement families $D_y$ that
appear in the finite lemma, now dense because the index set is infinite.

\subsection{Novelty is not genericity}\label{ssec:converse}

The finite lemma suggests a converse: perhaps a real that is new and splits every infinite coded
set is already generic. It is not.

\begin{proposition}[The converse fails; \condg, under \textbf{J1}]\label{prop:converse}
There is a total $x\colon\omega\to\{0,1\}$ such that $x\notin M$ and $x$ splits every infinite
$A\in\mathcal B_M$, but the restriction filter of $x$ is not generic over $M$.
\end{proposition}

\begin{proof}
By Lemma~\ref{lem:RS} take a generic $G$ and let $r=x_G$. Define the \emph{duplicated} real
$x(2n)=x(2n+1)=r(n)$. If $x$ were in $M$, so would be $r(n)=x(2n)$; hence $x\notin M$. Let $A\in M$
be infinite, and let $B=\{\lfloor a/2\rfloor:a\in A\}$. Then $B\in M$, and $B$ is infinite, since the
preimage of a finite $B$ would have at most $2|B|$ elements. By Theorem~\ref{thm:novelty}(2), $r$
splits $B$: there are $b,b'\in B$ with $r(b)\ne r(b')$. Choosing $a,a'\in A$ with $\lfloor
a/2\rfloor=b$ and $\lfloor a'/2\rfloor=b'$ gives $x(a)=r(b)\ne r(b')=x(a')$, so $x$ splits $A$.
Now consider
\[
  D^\ast=\{p:\ \text{for some }n,\ 2n,2n+1\in\dom p\text{ and }p(2n)\ne p(2n+1)\}.
\]
$D^\ast$ is defined without parameters, so it lies in $M$, and it is dense: above any condition,
choose a pair $2n,2n+1$ beyond its domain and assign $0$ and $1$. No restriction of $x$ lies in
$D^\ast$, so the restriction filter of $x$ misses a dense set in $M$.
\end{proof}

Genericity therefore asks for more than novelty and splitting. The exact requirement is meeting
\emph{every} dense set coded in $M$, and it has a clean topological form. Give $2^\omega$ the
product topology, whose basic open sets are the cylinders $[p]=\{x: p\subseteq x\}$. A \emph{tree} is a
set $T$ of conditions closed under restriction. A closed set $F\subseteq2^\omega$ is \emph{coded in
$M$} if $F=[T]$ for a tree $T\in M$, where $x\in[T]$ iff every finite restriction of $x$ lies in
$T$. The code is in $M$; its branches are computed outside $M$ and may include reals that are not
in $M$. Because the tree is finitely branching, emptiness of $[p]\cap[T]$ is decided at a finite
level:
\[
  [p]\cap[T]=\varnothing\iff\exists n\ \bigl(\dom p\subseteq n\ \wedge\ \forall s\in2^n\
  (p\subseteq s\Rightarrow s\notin T)\bigr).
\]
This criterion is arithmetic in $p$ and $T$, hence absolute for the transitive model $M$, and so
are the statements that $[T]$ is nowhere dense and that a family of conditions in $M$ is dense
(since $M$ contains exactly the actual finite conditions).

\begin{theorem}[Exact recognition of genericity; \condg, under \textbf{J1}]\label{thm:recognition}
For a total $x\colon\omega\to\{0,1\}$, the restriction filter $G_x=\{p:p\subseteq x\}$ is generic
over $M$ if and only if $x$ lies in no closed nowhere dense subset of $2^\omega$ coded in $M$.
\end{theorem}

\begin{proof}
($\Rightarrow$) Let $F=[T]$ be closed nowhere dense with $T\in M$. By the arithmetic criterion
above and Separation in $M$, the conditions $p$ with $[p]\cap F=\varnothing$ form a family in $M$,
and it is dense because $F$ is nowhere dense. A
generic $G_x$ meets it, so some $p\subseteq x$ has $[p]\cap F=\varnothing$, and $x\notin F$.
($\Leftarrow$) Let $D\in M$ be dense. Let $T_D$ be the set of conditions none of whose restrictions
lies in $D$; it is a tree, and it belongs to $M$ by Separation from $D$; and its branch set $F_D=[T_D]$ is the complement of the open set
$\bigcup_{p\in D}[p]$. Density of $D$ makes that open set dense, so $F_D$ is closed nowhere dense
and coded in $M$. If $x\notin F_D$, some restriction of $x$ lies in $D$, so $G_x$ meets $D$.
\end{proof}

The tests used in Theorem~\ref{thm:novelty} are special cases: avoiding the singleton $\{y\}$ for
$y\in M$ is meeting $D_y$, and avoiding the closed set of reals constant on an infinite $A\in M$ is
meeting $E_A$. They form a proper subfamily of all coded closed nowhere dense tests. The
duplicated real of Proposition~\ref{prop:converse} passes all of them and fails the test
belonging to $D^\ast$: it lies in the coded closed nowhere dense set of reals that are constant on
every pair $\{2n,2n+1\}$. Figure~\ref{fig:genericity} summarizes the relations. The density of
$D^\ast$, the fact that every duplicated real misses it, and the construction of the obstruction
trees $T_D$ with their nowhere density are verified in Lean (\S\ref{sec:mech}); the existence of $M$
and of generic reals is not.

\begin{figure}[tbp]
\centering
\begin{tikzpicture}[
    font=\footnotesize,
    reg/.style={draw, rounded corners=8pt, thick},
    dot/.style={circle, fill, inner sep=1.5pt}]
\draw[reg, fill=gray!6] (0,0) rectangle (14.2,6.4);
  \node[anchor=north west] at (0.15,6.35) {\textbf{all reals} $2^\omega$};
\draw[reg, fill=gray!22] (0.35,0.35) rectangle (2.75,5.3);
  \node[align=center] at (1.55,4.7) {reals in $M$\\(countable)};
  \node[dot] at (1.55,2.6) {};
  \node[anchor=north, align=center] at (1.55,2.45) {$y\in\mathcal B_M$\\fails its own\\test $D_y$};
\draw[reg, fill=orange!10] (3.05,0.35) rectangle (13.85,5.85);
  \node[anchor=north west] at (3.15,5.8) {\textbf{new}: outside $M$};
  \node[anchor=north east] at (13.75,5.8) {\emph{measure one}};
\draw[reg, fill=orange!22] (3.4,0.65) rectangle (13.5,5.1);
  \node[anchor=north west] at (3.5,5.05) {\textbf{new and splitting} every infinite
    $A\in\mathcal B_M$};
  \node[anchor=north east] at (13.4,5.05) {\emph{measure one}};
\draw[reg, fill=blue!12] (8.6,0.95) rectangle (13.2,4.2);
  \node[anchor=north west, align=left, text width=4.3cm] at (8.7,4.1)
       {\textbf{Cohen generic over $M$}\\[2pt]avoids every $M$-coded closed nowhere
        dense set (Thm.~\ref{thm:recognition})\\[4pt]\emph{comeagre, measure zero}};
\node[dot, red!70!black] at (4.2,2.4) {};
  \node[anchor=west, align=left, text=red!70!black] at (4.4,2.4)
       {duplicated real\\$x(2n)=x(2n+1)=r(n)$:\\new and splitting,\\but fails $D^\ast$
        (Prop.~\ref{prop:converse})};
\end{tikzpicture}
\caption{Where Cohen reals sit. Every generic real is new and splits every infinite coded set
(Theorem~\ref{thm:novelty}), so the blue region lies inside both orange ones. The converse
fails: the duplicated real (red) is new and splitting but misses the dense set $D^\ast$
(Proposition~\ref{prop:converse}). Genericity is characterized exactly by avoiding all coded closed
nowhere dense sets (Theorem~\ref{thm:recognition}). In measure the orange regions are full and the
blue region is null; in category the blue region is comeagre (Proposition~\ref{prop:category}).
Region sizes are schematic.}
\label{fig:genericity}
\end{figure}
 
\subsection{Strict extension of the coded predicates}\label{ssec:strict}

In the finite bridge, ``definable'' meant constant on the blocks of a partition. That picture
does not transfer literally to $\omega$.

\begin{proposition}[The coded predicates are not a partition lens; \inhouse, last sentence under \textbf{J1}]\label{prop:nolens}
Let $f\colon\omega\to X$ be any map. If every singleton predicate $\{k\}$ is constant on the fibers
of $f$, then $f$ is injective and \emph{every} predicate on $\omega$ is constant on its fibers. In
particular, under \textbf{J1}, $\mathcal B_M$ is not the family of predicates visible to any map on
$\omega$.
\end{proposition}

\begin{proof}
If $f(m)=f(n)$, the predicate $\{m\}$ takes the same value at $m$ and $n$, so $m=n$. An injective
map has singleton fibers, on which every predicate is constant. Under \textbf{J1},
$\mathcal B_M$ contains every singleton (each finite subset of $\omega$ is in $M$) but, being
countable, not every predicate.
\end{proof}

The right infinitary counterpart of the finite visible predicates is therefore the Boolean algebra
$\mathcal B_M$ itself, closed under the unions that $M$ can code but not under arbitrary external
countable unions. For a predicate $x$, let $\mathcal B_M[x]$ denote the Boolean algebra generated by
$\mathcal B_M\cup\{x\}$.

\begin{theorem}[Adjoining a generic is a strict extension; \condg]\label{thm:strict}
\begin{enumerate}\itemsep0pt
\item \inhouse. For any predicate $x$ on $\omega$, $\mathcal B_M\subseteq\mathcal B_M[x]$, and the
inclusion is strict iff $x\notin\mathcal B_M$.
\item Under \textbf{J1}, if $G$ is generic over $M$, then $\mathcal B_M\subsetneq\mathcal
B_M[x_G]$. Under \textbf{J2} in addition, $\mathcal B_M[x_G]\subseteq\powset(\omega)\cap M[G]$.
\end{enumerate}
\end{theorem}

\begin{proof}
(1) If $x\in\mathcal B_M$, every Boolean combination of $x$ with elements of $\mathcal B_M$ is
already in $\mathcal B_M$, so the two algebras coincide. If $x\notin\mathcal B_M$, then $x$ itself
is a new element. (2) The first claim is (1) together with Theorem~\ref{thm:novelty}(1). For the
second, $M[G]$ contains $x_G$ and every element of $\mathcal B_M$, and it is closed under finite
Boolean operations.
\end{proof}

The extension is strict on the same domain $\omega$: no new points are added, but the algebra of
expressible predicates grows. The framework's own formulation of strict extension is a failure of
factorization \citep{SBTCantorShell}, and Cohen forcing also satisfies that formulation, in a
deliberately coarse form.

\begin{theorem}[Non-factorization; \condg, under \textbf{J1}]\label{thm:nonfactor}
Let $\Gamma_M$ be the set of $\Cgen$-generic filters over $M$, and define $\pi_0(G):=M$ and
$\pi_1(G):=x_G$. Then $\pi_1$ does not factor through $\pi_0$: there is no $\varphi$ with
$\pi_1=\varphi\circ\pi_0$.
\end{theorem}

\begin{proof}
By Lemma~\ref{lem:RS} take generics $G^0\ni(0\mapsto0)$ and $G^1\ni(0\mapsto1)$. Then
$\pi_0(G^0)=\pi_0(G^1)$, while $x_{G^0}(0)=0\ne1=x_{G^1}(0)$. A map through the constant $\pi_0$
would be constant.
\end{proof}

Theorem~\ref{thm:nonfactor} uses a constant ground readout and so says little by itself; it does not
show that every observation available in $M$ fails to distinguish generics. The substantive
statement is Theorem~\ref{thm:strict}. Note also that strict extension of the predicate algebra is
not inclusion of first-order theories: $M$ and $M[G]$ can disagree about sentences of the old
language. If $M$ satisfies $V=L$, for example, then by \textbf{J2} $M[G]$ does not: $M[G]$ has the
same ordinals and hence the same constructible sets as $M$, and $x_G\notin M$.

\begin{corollary}[Cohen extension as a package change; \condg]\label{cor:strictcohen}
Adjoining a Cohen-generic real strictly enlarges the ground model's algebra of coded predicates
(Theorem~\ref{thm:strict}) by a predicate that no Boolean combination of ground predicates
produces; by \textbf{J2}, $M[G]$ is the model in which it lives. In the framework's terms this is a
package change \citep{SBTIncompleteness,SBTCantorShell}: new content obtained by adjoining
something the fixed package cannot express, the infinitary counterpart of the strict refinement
$(f,h)$ of a finite lens.
\end{corollary}

\subsection{Category, measure, and the counting lemma}\label{ssec:category}

The finite lemma is a statement about probability, while genericity is a statement about
category. Over $\omega$ the two come apart sharply.

\begin{proposition}[Cohen reals are comeagre but null; \condg, under \textbf{J1}]\label{prop:category}
Give $2^\omega$ the fair-coin measure. The reals $x$ with $G_x$ generic over $M$ form a comeagre set
of measure zero. By contrast, the reals outside $M$, and the reals that split every infinite
$A\in\mathcal B_M$, each form a set of measure one.
\end{proposition}

\begin{proof}
Genericity means lying in each of the countably many dense open sets $\bigcup_{p\in D}[p]$, for $D$
dense in $M$; their intersection is comeagre. For measure zero, enumerate the finite binary strings
as $s_0,s_1,\dots$ by a fixed arithmetic enumeration, and for $k,i\in\omega$ let $t_{k,i}$ be $s_i$
padded with zeros to length $\max(|s_i|,k+i+1)$. The open set $U_k=\bigcup_i[t_{k,i}]$ is dense and
has measure at most $\sum_i2^{-(k+i+1)}=2^{-k}$. The conditions extending some $t_{k,i}$ form a
dense family defined in $M$ without choice, so every generic real lies in every $U_k$, and the
generic reals form a null set. On the other side, $M$ contains only countably many reals, and for
each infinite $A$ the reals constant on $A$ form a null set; a countable union of null sets is
null.
\end{proof}

So the finite probability $1-2^{-(N-K)}\to1$ of meeting every disagreement family, which the
diagnostics of \S\ref{sec:mech} display, is the finite counterpart of novelty, which has measure
one, and not of full genericity, which has measure zero. Measure and category are dual notions of
largeness that disagree here, as they classically can \citep{Oxtoby1980}.

\begin{table}[tbp]\centering
\small
\begin{tabular}{@{}p{0.44\textwidth}p{0.50\textwidth}@{}}
\toprule
Finite bridge (\S\ref{ssec:finitebridge}) & Cohen forcing over a countable $M$ (\S\S\ref{ssec:vanish}--\ref{ssec:category}) \\
\midrule
finite index $Z$, $|Z|=N$ & infinite index $\omega$ \\
visible predicates: $2^K$ of $2^N$ & coded predicates $\mathcal B_M$: countably many of $2^{\aleph_0}$; a Boolean algebra, not a partition lens (Prop.~\ref{prop:nolens}) \\
predicate $h$, restriction filter $G_h$ & real $x_G$, generic filter $G=G_{x_G}$ \\
$D_g$ with defect $F(D_g)=\{g\}$ & $D_y$ dense for every $y$ (Cor.~\ref{cor:infinitecontrast}) \\
$E_x$ with defect at constant blocks & $E_A$ dense for every infinite $A$ \\
$h$ meets all $D_g$ iff $h$ is not definable & generic $\Rightarrow$ new and splitting (Thm.~\ref{thm:novelty}); not conversely (Prop.~\ref{prop:converse}) \\
every total $h$ meets every dense set (Rem.~\ref{rem:finitegeneric}) & generic $\Leftrightarrow$ avoids every coded closed nowhere dense set (Thm.~\ref{thm:recognition}) \\
$\Pr(\text{meets all }D_g)=1-2^{-(N-K)}$ & new reals: measure one; generic reals: comeagre, measure zero (Prop.~\ref{prop:category}) \\
strict refinement $(f,h)$ when $h$ is not definable & strict extension $\mathcal B_M\subsetneq\mathcal B_M[x_G]$ (Thm.~\ref{thm:strict}) \\
package change at $\omega$ (Thm.~\ref{thm:nonattain}) & generic adjunction (Cor.~\ref{cor:strictcohen}) \\
\bottomrule
\end{tabular}
\caption{The finite bridge and Cohen forcing side by side. Each row is backed by the indicated
result; rows 6--8 record where the finite picture does not transfer.}\label{tab:correspondence}
\end{table}

\subsection{The continuum hypothesis}\label{sec:ch}

\begin{theorem}[CH is not settled by the ground; \modelg]\label{thm:ch}
\begin{enumerate}[label=\textup{(\alph*)}]
\item If $\mathsf{ZF}$ is consistent, so is $\mathsf{ZFC}+\mathsf{CH}$
\citep{Godel1938,Godel1940}, and so is $\mathsf{ZFC}+\neg\mathsf{CH}$ \citep{Cohen1963}.
\item Assume, as a separate hypothesis, a countable transitive model $M$ of
$\mathsf{ZFC}+\mathsf{GCH}$. Adding $\aleph_2^M$ Cohen reals, with the poset
$\Fn(\aleph_2^M\times\omega,2)$, gives a generic extension $M[G]$ with the same ordinals and,
by the countable chain condition, the same cardinals as $M$; it has at least $\aleph_2$ distinct
reals, so $M\models\mathsf{CH}$ while $M[G]\models\neg\mathsf{CH}$
\citep{Kunen2011}.
\end{enumerate}
\end{theorem}

Part (a) is a pair of relative consistency results; it produces no countable transitive model and
places no two models on a common domain. Part (b) supplies such a picture under its own
hypothesis: the ground and the extension share their ordinals, though not their sets. Adding a
single Cohen real, as in the previous subsections, does not change the truth value of
$\mathsf{CH}$; many are needed. In the language of this paper, the predicates a ground model can
code do not determine the size of the continuum: the answer depends on which extension is taken.

\begin{remark}[The set layer is not a unique package]\label{rem:notunique}
Theorem~\ref{thm:ch} is the positive form of one of our non-claims: we do not claim that
$\mathsf{ZFC}$ is the unique or canonical set theory. Reading such disagreements as distinct and
equally legitimate universes is the set-theoretic multiverse view of \citet{Hamkins2012}; the
universe view (e.g.\ \citealp{Godel1947}) holds that one of them is correct. We construct no
model, reprove no independence result, and take no position between these views. On the reading
offered here, independence is what one expects of a foundation that grows by lawful extension.
\end{remark}
 \section{Guards: a countermodel atlas}\label{sec:atlas}

The claims of this paper are easy to overstate by a single word, and the overstatements are the
ones a reader is most tempted to supply: ``extensionality is forced,'' ``Foundation is free,''
``Choice falls out,'' ``a new real is a generic real.'' Each is false, and each is blocked by a
specific countermodel. A free/purchased reading is meaningful only if, for each claim that
something is free, one can exhibit the alternative under which it fails. Each entry below is
\guardg{}; none proves a positive theorem, and most of the witnesses already appear above.

\begin{description}[leftmargin=1.2em, itemsep=4pt]

\item[\textbf{CM-1}: Extensionality is a choice of lens.] The membership lens is not injective on
presentations (Lemma~\ref{lem:L2}): the two-node and three-node presentations of
$\{\varnothing\}$, the latter with a duplicated childless node, have the same profile and are
identified. A different instrument keeps them apart. For example, the lens $f'(P)=(f(P),|V(P)|)$,
which also records the number of nodes, distinguishes the two presentations, and the identity map
serves as its packaging; so does the lens that records the full pointed-graph isomorphism type.
This blocks the reading that the framework \emph{forces} extensionality, while agreeing that the
canonical instrument satisfies it. (Anti-foundation is not an example here: Aczel's theory keeps
Extensionality and also identifies the duplicated childless nodes; what it changes is Foundation,
\textbf{CM-2}.)

\item[\textbf{CM-2}: Foundation needs staging.] The self-loop $x=\{x\}$ and the two-cycle
$x=\{y\}$, $y=\{x\}$ are coherent finite membership graphs that admit \emph{no} strictly
decreasing audit (Theorem~\ref{thm:tfae}). They are excluded only because the present instrument
was staged acyclic. This blocks ``Foundation is free without staging''; anti-foundation
\citep{Aczel1988} is the classical theory that admits such graphs.

\item[\textbf{CM-3}: Infinity does not emerge from iteration.] $\HF$ is the witness: no finitely
generated rule started inside $\HF$ leaves it (Theorem~\ref{thm:nonattain}). An infinite completed
object may be \emph{admitted}, but it is not \emph{reached}.

\item[\textbf{CM-4}: Sections are not free infinitarily.] If $\mathsf{ZF}$ is consistent, there
are models of $\mathsf{ZF}$ in which some surjection has no section; Cohen's symmetric models are
examples \citep{Cohen1966,Jech1973}. This blocks ``\textbf{C-SEL} is derivable in $\mathsf{ZF}$.''
Finite sections, and sections of lenses that come with an explicit section or a well-ordered
source, remain free (Theorem~\ref{thm:choice}).

\item[\textbf{CM-5}: Replacement is not supplied by the other commitments.] In ambient
$\mathsf{ZFC}$, $V_{\omega+\omega}$ satisfies (S), \textbf{C-INF}, \textbf{C-POW} and
\textbf{C-SEL} but not Replacement: the image of $n\mapsto\omega+n$ is not an element
(Proposition~\ref{prop:nonredundancy}(ii)).

\item[\textbf{CM-6}: Not every predicate extension is strict.] If $g$ is already definable from
$f$, the refined lens $(f,g)$ induces the same partition as $f$, and $G_g$ misses $D_g$
(Lemma~\ref{lem:meeting}). Likewise, adjoining a real that is already in $M$ leaves
$\mathcal B_M$ unchanged (Theorem~\ref{thm:strict}(1)). Strictness needs non-definability, which
the counting lemma (Theorem~\ref{thm:counting}) shows is typical but not universal.

\item[\textbf{CM-7}: Power set is not supplied by the other commitments.] In ambient
$\mathsf{ZFC}$, $H_{\omega_1}$ satisfies (S), \textbf{C-INF}, \textbf{C-REPL} and \textbf{C-SEL}
but has no power set of $\omega$ (Proposition~\ref{prop:nonredundancy}(iii)).

\item[\textbf{CM-8}: A new real need not be generic.] The duplicated real $x(2n)=x(2n+1)=r(n)$,
with $r$ generic, is outside $M$ and splits every infinite coded set, yet misses the dense set
$D^\ast$ (Proposition~\ref{prop:converse}). This blocks ``Cohen genericity is
non-definability''; only the forward implication holds (Theorem~\ref{thm:novelty}), and the exact
test is Theorem~\ref{thm:recognition}.

\end{description}

\noindent By provenance, \textbf{CM-1}, \textbf{CM-2} and the finite part of \textbf{CM-6} are
finite witnesses; \textbf{CM-3} is the in-house non-attainment theorem, a statement about the
infinite family $\HF$ made in the ambient metatheory; the ground-model part of \textbf{CM-6} and
\textbf{CM-8} are proved under \textbf{J1}; and \textbf{CM-4}, \textbf{CM-5} and \textbf{CM-7}
are classical, recalled in relative-consistency form or inside ambient $\mathsf{ZFC}$. Together they delimit what the
free/purchased reading does, and does not, assert.
 \section{Mechanization and computational evidence}\label{sec:mech}

Selected finite and structural cores of the paper are verified in Lean, and the counting results
are checked against reproducible numerical exhibits. This section states what is covered, what is
not, and what the axiom audit shows.

\subsection{A Lean development without external libraries}

The development uses Lean~4 \citep{deMouraUllrich2021}, version 4.28.0, and imports no external
mathematical library---in particular not \textsf{mathlib} \citep{Mathlib2020}, whose
\texttt{ZFSet} construction it conceptually mirrors (Remark~\ref{rem:mathlib}). Membership trees,
extensional equality, the Cohen posets, rank and the idempotent-split structure are all defined
from first principles. Eight modules carry the verified core:
\begin{description}[leftmargin=1.2em, itemsep=2pt]
\item[\texttt{Packaging}.] Idempotents split and have fixed points equal to their image, with a
finite collapse witness. This is the abstract pattern of Theorem~\ref{thm:p5}; the construction
of canonical presentations from graphs is not mechanized.
\item[\texttt{HFClosure}.] Finite membership trees with recursive extensional equality:
extensional equality holds iff the members agree; pairing, union and filtering respect it;
finite collection is derived; every family that respects equality and is closed under empty set,
pairing and union contains every tree (the core of Theorem~\ref{thm:hfclosure}); and no tree is
inductive.
\item[\texttt{HFSets}.] The quotient of finite trees by extensional equality, formed without
choosing representatives: Extensionality, Empty set, Pairing, Union, Separation for decidable
(Boolean-valued) predicates, rank decrease along membership, no self-membership, and no inductive
set. Separation is verified for Boolean predicates on the quotient, not for the full formula
schema of Theorem~\ref{thm:hfsat}.
\item[\texttt{RankAudit}.] A strictly decreasing audit excludes cycles and yields minimal
elements of nonempty node lists; the self-loop and two-cycle of \textbf{CM-2} are cycles and so
admit no audit. The implication from acyclicity to an audit is left on paper.
\item[\texttt{FiniteCohen}.] The meeting equivalences of Lemma~\ref{lem:meeting}, the density
defects of Theorem~\ref{thm:densitydefect}, and Remark~\ref{rem:finitegeneric}: in the finite poset
every total predicate meets every dense family, and $D_g$ is not dense.
\item[\texttt{CohenBoundary}.] Conditions on $\mathbb N$ with finite support: $D_g$ is dense and is
met exactly when a disagreement witness exists; $E_A$ is dense for unbounded $A$
(Corollary~\ref{cor:infinitecontrast}); $D^\ast$ is dense and every duplicated real misses it
(Proposition~\ref{prop:converse}); visible singletons force injectivity
(Proposition~\ref{prop:nolens}); and adjoining $x$ to a Boolean-closed family is strict iff $x$ is
new (Theorem~\ref{thm:strict}(1)).
\item[\texttt{CohenRecognition}.] The obstruction trees of Theorem~\ref{thm:recognition}: for a dense
family $D$, the tree $T_D$ is downward closed and nowhere dense, and a real meets $D$ iff it is not
a branch of $T_D$. Genericity is formalized relative to an explicitly supplied register of coded
families. It is equivalent to avoiding the obstructions of an indexed list of families provided
every listed family is coded and dense and every coded dense family is listed. It implies novelty
with respect to a family $\mathcal B$ of ground predicates provided the disagreement family $D_g$ of
every $g\in\mathcal B$ is registered, and strict extension provided in addition that $\mathcal B$ is
Boolean-closed. Once $D^\ast$ is registered, duplicated reals are excluded. With an empty register
every real would count as generic, so these premises are essential.
\item[\texttt{LiftChoice}.] Choice-free sections of surjections with finite source, and the
identity between sections and prototype families (Theorem~\ref{thm:choice}, finite part).
\end{description}

Several things are deliberately left on paper: the isomorphism between rooted extensional graphs
and canonical presentations; satisfaction of the full Separation and Replacement schemas in
$(\HF,\in)$; non-attainment for general finitely generated rules; the rank-erasure translation;
the counting lemma; the $\mathsf{ZF}$ equivalences of Choice; and everything about countable
transitive models, generic filters over them and $M[G]$. In particular, \texttt{CohenRecognition}
takes the register of coded dense families as an input; that this register consists of exactly
the dense sets belonging to a model $M$ is the mathematical bridge argued in \S\ref{sec:forcing},
not a Lean theorem.

\begin{proposition}[Axiom base; \inhouse]\label{prop:axioms}
Every declaration of the Lean development, including compiler-generated auxiliary declarations,
depends on no axiom other than propositional extensionality and quotient soundness
\citep{Carneiro2019}. In particular, \texttt{Classical.choice}, the third axiom of Lean's standard
classical setup, is used nowhere, and no declaration uses an unfinished proof or a project axiom.
\end{proposition}

The proposition is checked mechanically, not by inspection: an audit file, built as a default
target of every \texttt{lake build}, selects declarations by their defining module, so that private auxiliary declarations are
included, and traverses all 371 of them (194 theorems, 9 private declarations) with Lean's
transitive axiom collector; it fails the build if any other axiom appears. As negative controls,
temporary files that inject a project axiom, an unfinished proof, or a use of
\texttt{Classical.choice} are each rejected by the same check. Where a proof needs classical reasoning, as in the
converse direction of the recognition theorem, the relevant instance of excluded middle is stated
as an explicit hypothesis rather than imported through \texttt{Classical.choice}.

This matters for the argument of \S\ref{sec:purchases}. The paper prices Choice as the cost of
sections for arbitrary surjections; a verification of the finite, free case that relied on Choice
would be circular. The mechanized finite sections set that price to zero without Choice, while the
equivalence with $\mathsf{AC}$ for arbitrary sets is proved on paper in $\mathsf{ZF}$.

\subsection{Reproducible numerical exhibits}

Three deterministic, seeded scripts check the quantitative predictions on finite data. Each
asserts a prediction that could fail and reports the residual; every reported number is
regenerated from a stored run rather than transcribed. The exhibits illustrate the theorems and
check the encodings; they prove nothing beyond the bounded cases they examine.

\paragraph{The $\HF$ tower.} The sets of rank at most $n$ form $V_{n+1}$, so the enumeration gives
$|V_1|,\dots,|V_5|=1,2,4,16,65536$. Over all $65{,}536$ elements of $V_5$ the exhibit checks
that Ackermann coding and decoding round-trip, that rank strictly decreases along all
$524{,}288$ membership edges (Theorem~\ref{thm:tfae}), and that no element is inductive
(Theorem~\ref{thm:hfsat}). It also computes the power sets of all $16$ elements of $V_4$ and the
images of all $625$ maps from those sets into $V_3$, and samples pairing, union, separation and
finite collection.

\paragraph{Definability rarity.} For nine lens configurations with $N-K$ between $0$ and $20$, the
fraction of $2\times10^5$ sampled predicates that are definable matches $2^{-(N-K)}$ within a
statistical gate wherever the expected count is large enough to test, and the fraction that
splits every non-singleton block matches the exact product formula of Theorem~\ref{thm:counting}
(Figure~\ref{fig:diagnostics}).

\paragraph{Finite Cohen bridge.} An exhaustive enumeration over the empty lens and all
surjective labelled lenses on at most five points ($634$ lenses) confirms the meeting equivalences
and density defects of Lemma~\ref{lem:meeting} and Theorem~\ref{thm:densitydefect} on every
instance, together with the counts $2^K$ and the splitting product. A scaling experiment confirms
that the probability of meeting every disagreement family is $1-2^{-(N-K)}$. As explained in
\S\ref{ssec:category}, this is a finite picture of novelty, not of Cohen genericity.

\begin{figure}[tbp]
\centering
\begin{tikzpicture}
\begin{axis}[
    name=left,
    width=0.47\textwidth, height=0.36\textwidth,
    ymode=log, xmin=-1, xmax=21, ymin=3e-7, ymax=3,
    xlabel={$N-K$}, ylabel={fraction definable from $f$},
    label style={font=\small}, tick label style={font=\scriptsize},
    legend style={font=\scriptsize, at={(0.97,0.97)}, anchor=north east, draw=none,
                  fill=white, fill opacity=0.8, text opacity=1},
    title={\small (a) definability rarity}, title style={yshift=-4pt}]
  \addplot[domain=0:20, samples=41, thick, blue!60!black] {2^(-x)};
  \addlegendentry{exact $2^{-(N-K)}$}
  \addplot[only marks, mark=x, mark size=3pt, thick, orange!85!black,
           unbounded coords=discard]
    table[x=nk, y=emp] {figures/data/rarity.dat};
  \addlegendentry{empirical ($2\times10^5$ samples)}
\end{axis}
\begin{axis}[
    at={(left.east)}, anchor=west, xshift=1.6cm,
    width=0.40\textwidth, height=0.36\textwidth,
    xmin=0, xmax=1.02, ymin=0, ymax=1.02,
    xlabel={exact $\prod_{|B|\ge2}(1-2^{1-|B|})$}, ylabel={empirical},
    label style={font=\small}, tick label style={font=\scriptsize},
    title={\small (b) splitting every block}, title style={yshift=-4pt}]
  \addplot[domain=0:1, gray, dashed] {x};
  \addplot[only marks, mark=*, mark size=1.8pt, blue!60!black]
    table[x=split_exact, y=split_emp] {figures/data/rarity.dat};
\end{axis}
\end{tikzpicture}
\caption{Finite counting checks (Theorem~\ref{thm:counting}) on the nine lens configurations of
the definability-rarity exhibit. (a) Fraction of sampled predicates that are constant on every
block, against the exact $2^{-(N-K)}$. The two configurations with $N-K=8$ have different block
sizes but the same theoretical value, and both empirical values pass the statistical gate. At
$N-K=20$ the expected number of definable samples is about $0.2$ and none was observed, so no point
is drawn in (a); that definability measurement is recorded as an observation, not a statistical
test. (b)
Fraction of predicates that split every block of size at least two, against the exact product
formula. The plotted values are read from the stored diagnostic runs, not re-sampled.}
\label{fig:diagnostics}
\end{figure}
  \section{Scope and non-claims}\label{sec:scope}

The results of this paper differ sharply in strength, and their value depends on keeping them
apart. This section collects the boundaries.

\paragraph{Proved here (\inhouse).} The finite tier: the packaging quotient and its
identification with $\HF$ (Theorem~\ref{thm:p5}); the closure theorem and the satisfaction of
every $\mathsf{ZFC}$ axiom except Infinity by $(\HF,\in)$ (Theorems~\ref{thm:hfclosure},
\ref{thm:hfsat}); the rank-audit equivalence (Theorem~\ref{thm:tfae}); non-attainment of
inductive sets by finitely generated rules (Theorem~\ref{thm:nonattain}); the rank-erasure
translation (Theorem~\ref{thm:translation}); the finite Cohen bridge, its density defects and the
counting lemma (Lemma~\ref{lem:meeting}, Theorems~\ref{thm:densitydefect} and \ref{thm:counting});
the infinite-index density facts (Corollary~\ref{cor:infinitecontrast}); the fact that a map on
$\omega$ under which all singletons are visible is injective (Proposition~\ref{prop:nolens}, whose
consequence for $\mathcal B_M$ uses \textbf{J1}); the algebraic strictness
criterion (Theorem~\ref{thm:strict}(1)); and the finite sections together with the standard
$\mathsf{ZF}$ equivalences used to price Choice (Theorem~\ref{thm:choice}). Global statements about
$\HF$ are made in the ambient set-theoretic metatheory. The Lean development verifies selected
cores of these results, listed in \S\ref{sec:mech}.

\paragraph{Proved under named hypotheses (\condg).} The shadow theorem
(Theorem~\ref{thm:shadow}) assumes the structural premise (S) and the commitments
\textbf{C-INF}, \textbf{C-POW}, \textbf{C-REPL} and, for Choice, \textbf{C-SEL}. The forcing
results (Theorems~\ref{thm:novelty}, \ref{thm:recognition}, \ref{thm:strict}(2),
\ref{thm:nonfactor}, Propositions~\ref{prop:converse} and \ref{prop:category}) assume \textbf{J1};
\textbf{J2} is used only to identify the extension $M[G]$. The non-redundancy witnesses
(Proposition~\ref{prop:nonredundancy}(ii)--(iv)) are classical and are used in ambient
$\mathsf{ZFC}$ or in relative-consistency form. Removing these hypotheses removes precisely the
conditional results and leaves the free part standing.

\paragraph{Statements about specific models (\modelg).} The continuum-hypothesis statements
(Theorem~\ref{thm:ch}) are cited classical facts; part (b) needs its own hypothesis of a countable
transitive model of $\mathsf{ZFC}+\mathsf{GCH}$.

\paragraph{What we do \emph{not} claim.}
\begin{enumerate}\itemsep1pt
  \item We do not derive $\mathsf{ZF}$ or $\mathsf{ZFC}$ from finite principles. The infinitary
  content enters through the structural premise and the four commitments, which are hypotheses.
  \item We prove no consistency statement beyond the classical results cited. Nothing here
  implies the existence of a countable transitive model, of an inaccessible cardinal, or the
  consistency of $\mathsf{ZFC}$.
  \item We claim no new independence result. The forcing section reorganizes classical facts and
  adds finite and structural statements about them.
  \item We do not claim that Cohen genericity is the same as non-definability from the ground
  model. Genericity implies novelty and splitting; the converse is false
  (Proposition~\ref{prop:converse}), and genericity is characterized only by the full family of
  coded closed nowhere dense tests (Theorem~\ref{thm:recognition}).
  \item We do not claim that the finite probabilities measure Cohen genericity. They describe
  novelty; Cohen reals have measure zero (Proposition~\ref{prop:category}).
  \item We do not claim that $\mathsf{ZFC}$ is the unique or most natural set theory; the
  continuum-hypothesis statements point the other way.
  \item We make no metaphysical claim about a universe of sets. ``Set'' is used as a package and
  ``$\in$'' as a lens; the readings are structural.
  \item The reading of all this as an account of where the axioms come from
  (\S\ref{sec:discussion}) is interpretation (\metag). It comes after the theorems and is never a
  premise for any of them.
\end{enumerate}

\paragraph{On the G\"odel--Cohen comparison.} The recurring reading of incompleteness and
independence as the saturation and extension faces of one discipline
(Remarks~\ref{rem:saturation} and \ref{rem:notunique}, \S\ref{sec:discussion}) is structural
framing, marked \metag. Each half is backed by theorems---saturation by
Theorem~\ref{thm:nonattain}, extension by Theorems~\ref{thm:novelty} and \ref{thm:strict}---but
the unification of the two is a way of seeing, not a further theorem.
 \section{Where the axioms come from}\label{sec:discussion}

This closing section is interpretation (\metag). It contains no theorem; each paragraph draws on a
result above, but the picture they compose is a reading, not a further result.

\subsection{Two limitative theorems, one pattern}

The foundations of mathematics carry two famous limits. G\"odel's theorems \citep{Godel1931} say
that a consistent, effectively axiomatized theory containing enough arithmetic does not prove
every arithmetical truth and, under the usual conditions on provability, does not prove its own
consistency. Cohen's theorems \citep{Cohen1963} say that natural questions, the continuum
hypothesis first among them, are not settled by the standard axioms. The two are usually filed
separately, as a limit on proof and a limit on decision, and usually as bad news.

The construction of this paper suggests filing them together. Its organizing patterns are that a
fixed rule, iterated, \emph{saturates}---it stays inside a closure it cannot leave---and that
content outside that closure comes only from \emph{extension}, the adjunction of something the
current package cannot express. Both patterns appear at the scale of
set theory. Saturation is Theorem~\ref{thm:nonattain}: the finite packaging rules never leave
$\HF$, which is why Infinity has to be purchased. Extension is the forcing analysis: a generic real
lies outside the ground model and strictly enlarges its algebra of coded predicates
(Theorems~\ref{thm:novelty} and \ref{thm:strict}).

On this reading G\"odel is the saturation face: a fixed, effectively given theory, run as long as
one likes, does not reach certain truths about itself. Cohen is the extension face: a ground model
is enlarged by adjoining a real it cannot code, and in suitable extensions old questions, such as
the size of the continuum, receive different answers. Two cautions keep the reading honest. The
G\"odel half needs all of its hypotheses: complete consistent theories exist, such as true
arithmetic, but they are not effectively axiomatized. And an independent \emph{sentence} is not the
same thing as a missing \emph{predicate}: Theorem~\ref{thm:strict} is about predicates on $\omega$,
and nothing here identifies the sentences a theory leaves undecided with the reals a model fails to
contain. The resemblance is one of shape---a closed apparatus run to saturation, and the lawful
extension beyond it---and it is offered as such.

\subsection{A foundation is an audited ledger}

If this is right, a foundation is neither a revealed truth nor an arbitrary convention, but a
\emph{ledger}. The distinction it draws, between questions settled inside a chosen framework and
the external question of which framework to adopt, is Carnap's \citep{Carnap1950}.
Table~\ref{tab:factorization} is that ledger written out for $\mathsf{ZFC}$. The free entries are
supplied by a modest packaging instrument and hold in $\HF$, with selected finite cores
machine-checked; at infinite stages they are carried by a structural premise that applies the same
clauses everywhere. The purchased entries are four named commitments, each shown to be
non-redundant, and the first boundary is sharp: $\HF$ satisfies everything except Infinity
(Proposition~\ref{prop:nonredundancy}(i)).

This is how $\mathsf{ZFC}$ can be a flat list and still be faithful to the staged construction:
over $\HF$ the staging is recoverable from membership alone (Theorem~\ref{thm:translation}). The
list is also not arbitrary, since its dependencies are fixed by what packaging supplies and what
must be bought. One can choose differently---a different lens (\textbf{CM-1}), no Foundation
(\textbf{CM-2}), no power-set totalization (\textbf{CM-7}), no selection principle
(\textbf{CM-4})---and the countermodel atlas is a catalogue of such alternative ledgers, in the
limit the plurality of set-theoretic universes discussed by \citet{Hamkins2012}. On this view
$\mathsf{ZFC}$ is one audited ledger among several, distinguished not by being true of an
antecedently given universe of objects \citep{Benacerraf1965} but by which commitments it buys.

\subsection{The open questions are expected}

On this reading the famous open questions lose some of their air of scandal. A foundation that
grows by lawful extension should have undecided questions: places where different extensions give
different answers. The continuum hypothesis is the standard example (\S\ref{sec:ch}). Over a
suitable ground model, one extension keeps it and another refutes it, and the ground's coded
predicates do not decide between them. Independence, read this way, is the expected signature of a
foundation that is still being extended, not a defect of one that was supposed to be finished.

\subsection{The answer}

We began with a list and a question: where does the list come from? Not from a pre-existing
universe of sets that the axioms approximately describe; we built no such universe and needed
none. Not from arbitrary convention either: the free entries are supplied by the chosen membership
instrument, the purchases are priced, and the order of dependence is not ours to choose. The list
comes from the place the framework locates stable objects in general: \emph{packaging under limited
access}---collecting multiplicities into units seen only through a lens---\emph{audited so that the
bookkeeping stays honest}, and \emph{extended when it saturates}. Set theory flattens that history
into a single membership relation; rank-erasure is the flattening; and forcing shows that the
history can always be continued.\footnote{The question ``where mathematics comes from'' is also
the title of \citet{LakoffNunez2000}, whose thesis concerns the cognitive origin of mathematical
ideas. That question is orthogonal to the structural account of where the \emph{axioms} of a
foundation come from offered here.}
 
\appendix
\section{Glossary of terms and hypotheses}\label{app:definitions}

This appendix collects terms used across several sections. It introduces no new mathematics; the
controlling definitions are the ones cited.

\subsection{Terms}\label{subsec:definitions-list}

\begin{description}[style=multiline,leftmargin=!,labelwidth=0.27\textwidth,labelsep=0.8em,itemsep=3pt]
\item[Theory package.] A carrier, lens, packaging map and audit $(Z,f,E,\mathcal A)$
(Convention~\ref{conv:package}), used here as an organizing instrument.
\item[Presentation.] A finite pointed acyclic graph $(V,\to,v)$ with $V\subseteq\mathbb N$ finite,
read as a membership presentation (Definition~\ref{def:presentation}). The set $Z$ of all
presentations is countably infinite.
\item[Membership-profile lens.] The map $f\colon Z\to\HF$ sending a presentation to the decoration
of its point (Definition~\ref{def:lens}).
\item[Packaging map.] $E(P)=C(f(P))$, which forgets the shape of a presentation and re-presents its
profile canonically (Definition~\ref{def:E}).
\item[Rank audit.] $\rho(P)=\rank(f(P))$, strictly decreasing along membership
(Definition~\ref{def:rank}, Lemma~\ref{lem:L3}).
\item[Finite packaging operator.] The operator $F$ on subfamilies of $\HF$ generated by the empty
set, pairing, union, subsets (gates) and finite images (transport) (Definition~\ref{def:F}).
\item[Finitely generated rule.] A completion rule adding only finite sets whose members come from
the current family and the transitive closures of its members (Definition~\ref{def:fg}).
\item[Rank-erasure.] Passing from the two-sorted staged structure to the one-sorted membership
structure (Definition~\ref{def:staged}).
\item[Coded predicates $\mathcal B_M$.] $\powset(\omega)\cap M$, a countable Boolean algebra of
predicates on $\omega$ (Definition~\ref{def:groundlens}); not the visible-predicate family of any
map on $\omega$ (Proposition~\ref{prop:nolens}). $\mathcal B_M[x]$ is the Boolean algebra generated
by $\mathcal B_M$ and $x$.
\item[Restriction filter.] For a total $x$, $G_x=\{p:p\subseteq x\}$.
\item[$D_g$, $E_A$, $D^\ast$.] Disagreement with $g$; splitting of $A$; disagreement within some pair
$\{2n,2n+1\}$ (\S\S\ref{ssec:finitebridge}, \ref{ssec:converse}).
\item[Coded closed set.] The branch set $[T]$ of a tree of conditions $T\in M$, computed outside
$M$ (\S\ref{ssec:converse}).
\item[Strict extension.] An enlargement of the algebra of expressible predicates on the same
domain, $\mathcal B_M\subsetneq\mathcal B_M[x]$ (Theorem~\ref{thm:strict}); in the framework's
coarser formulation, a failure of factorization (Theorem~\ref{thm:nonfactor}). It is not an
inclusion of first-order theories.
\item[Package change.] New content obtained by adjoining something the fixed package cannot produce
or express (Corollaries~\ref{cor:packagechange} and \ref{cor:strictcohen}).
\end{description}

\subsection{Hypotheses and regimes}\label{subsec:gate-predicates-list}

\begin{description}[style=multiline,leftmargin=!,labelwidth=0.27\textwidth,labelsep=0.8em,itemsep=3pt]
\item[\textbf{A\_FIN}.] The framework's finite-carrier regime. Here it applies to each
presentation and each element of $\HF$ separately: finite gates, finite transport, least-index
sections and finite diagnostics. It does not license any infinitary commitment.
\item[\textbf{(S)}.] The structural premise: a transitive admitted universe $U$ closed under empty
set, pairs, unions and relatively definable subsets (\S\ref{ssec:box}).
\item[\textbf{C-INF}.] An inductive set in $U$; priced by Theorem~\ref{thm:nonattain}.
\item[\textbf{C-POW}.] For each $x\in U$, an element of $U$ whose members are exactly the subsets of
$x$ lying in $U$; its finite form is Corollary~\ref{cor:finpow}.
\item[\textbf{C-REPL}.] Images in $U$ of sets in $U$ under maps definable in $U$; its finite form
is the transport clause.
\item[\textbf{C-SEL}.] Sections in $U$ for all surjections in $U$; equivalent to Choice in $U$ when
$U$ satisfies $\mathsf{ZF}$ (Corollary~\ref{cor:choice}).
\item[\textbf{J1}.] A countable transitive model $M$ of $\mathsf{ZFC}$, countable externally.
\item[\textbf{J2}.] The classical generic-extension machinery for Cohen forcing, used only to
identify $M[G]$.
\item[Diagnostic gates.] Pass/fail criteria in the computational exhibits. They test encodings and
numerical predictions on bounded cases and do not raise any statement to theorem grade.
\end{description}
 \section{Methods and reproducibility}\label{app:methods}

The paper rests on two reproducible support layers, a Lean development and deterministic
diagnostics. This appendix records how they are run and how they are used. No artifact is promoted
beyond the grade assigned in the body.

\paragraph{Specification.} The theorem statements, proof obligations, non-claims and grade
ceilings are recorded in \path{docs/spec/paper-contract.md} and
\path{docs/spec/theorem-inventory.md}. The mathematical review of all statements, including the
counterexample and characterization of \S\ref{ssec:converse} and the scope corrections they
entail, is recorded in \path{docs/review/mathematical-audit.txt}.

\paragraph{Lean.} The development lives in \path{formal/} and builds with \texttt{cd formal \&\&
lake build} on the pinned toolchain \texttt{leanprover/lean4:4.28.0}, with no external Lean
dependencies. The axiom audit \path{formal/AxiomAudit.lean} is a default build target, so an
ordinary build fails if any project declaration depends on an axiom other than
\textsf{propext} or \textsf{Quot.sound}.

\paragraph{Diagnostics.} The scripts \path{diagnostics/hf_slice.py},
\path{diagnostics/definability_rarity.py} and \path{diagnostics/finite_cohen.py} run deterministic
seeded computations under the pinned dependencies in \path{diagnostics/requirements.txt}. Each
writes a JSON record to \path{diagnostics/results/}, regenerates its tables in
\path{diagnostics/tables/}, and exits with a counterexample record if a check fails. Statistical
checks use binomial residual gates where the expected counts are large enough; cases with small
expected counts are recorded as observations without a pass criterion. The data for
Figure~\ref{fig:diagnostics} are extracted from the stored JSON records by
\path{paper/figures/data/make_fig_data.py}, without re-sampling.

\paragraph{Scope.} The Lean build checks selected finite and structural cores; it does not formalize
the imported infinitary set theory. The diagnostics check bounded cases and the agreement of their
tables with their records; they do not replace proofs. The classical forcing, cumulative-hierarchy,
natural-model, choice-equivalence and countermodel results remain external dependencies, cited
where used.

\paragraph{Availability.} All sources, the Lean development, the diagnostic scripts and their
generated artifacts are available at \url{https://github.com/ioannist/six-birds-zfc}. No external
datasets are used.
 \section{Results and their artifacts}\label{app:artifact-map}

Table~\ref{tab:artifact-map} maps the main results to the Lean declarations that verify their
cores (module \path{formal/SetStone/<Module>.lean}) and to the diagnostics that exhibit them. A
dash means that the result is proved on paper only; \S\ref{sec:mech} lists what each module does
and does not cover.

\begin{table}[htbp]
\centering
\footnotesize
\begin{tabularx}{\textwidth}{@{}>{\raggedright\arraybackslash}p{0.22\textwidth}>{\raggedright\arraybackslash}X>{\raggedright\arraybackslash}p{0.2\textwidth}@{}}
\toprule
Result & Lean module: key declarations & Diagnostic \\
\midrule
Thm.~\ref{thm:p5} (packaging) & \texttt{Packaging}: \texttt{fix\_iff\_image}, \texttt{r\_i\_id}, \texttt{i\_r\_eq\_e}, \texttt{collapseFin3\_idem} & --- \\
Lem.~\ref{lem:L2}, Thm.~\ref{thm:hfclosure} & \texttt{HFClosure}: \texttt{equiv\_iff\_mem}, \texttt{collect\_in\_closed}, \texttt{closed\_contains\_all}; \texttt{HFSets}: \texttt{ext}, \texttt{closed\_family\_all} & \path{hf_slice.py} \\
Thm.~\ref{thm:hfsat} (selected axioms) & \texttt{HFSets}: \texttt{mem\_pair}, \texttt{mem\_union}, \texttt{mem\_separate}, \texttt{not\_mem\_self}, \texttt{no\_inductive\_set}; \texttt{HFClosure}: \texttt{not\_inductiveHF} & \path{hf_slice.py} \\
Lem.~\ref{lem:L3}, Thm.~\ref{thm:tfae} & \texttt{HFSets}: \texttt{mem\_rank\_lt}; \texttt{RankAudit}: \texttt{potential\_excludes\_cycle}, \texttt{min\_of\_potential}, \texttt{self\_loop\_is\_cycle} & \path{hf_slice.py} \\
Thm.~\ref{thm:nonattain}, Thm.~\ref{thm:translation} & --- & --- \\
Thm.~\ref{thm:choice} (finite part) & \texttt{LiftChoice}: \texttt{finite\_section\_exists}, \texttt{section\_iff\_prototype} & --- \\
Lem.~\ref{lem:meeting}, Thm.~\ref{thm:densitydefect}, Rem.~\ref{rem:finitegeneric} & \texttt{FiniteCohen}: \texttt{meet\_D\_iff\_neq}, \texttt{meet\_E\_iff\_splits}, \texttt{density\_D\_failure\_iff\_total}, \texttt{density\_E\_failure\_iff\_constant}, \texttt{every\_total\_meets\_dense}, \texttt{disagreement\_not\_dense} & \path{finite_cohen.py} \\
Thm.~\ref{thm:counting} & --- & \path{definability_rarity.py}, \path{finite_cohen.py} \\
Cor.~\ref{cor:infinitecontrast} & \texttt{CohenBoundary}: \texttt{disagreement\_dense}, \texttt{disagreement\_meeting\_iff}, \texttt{infinite\_block\_splitting\_dense} & --- \\
Prop.~\ref{prop:converse} (finite core) & \texttt{CohenBoundary}: \texttt{pairMismatch\_dense}, \texttt{duplicate\_misses\_dense} & --- \\
Thm.~\ref{thm:recognition} (finite core) & \texttt{CohenRecognition}: \texttt{obstruction\_nowhereDense}, \texttt{meets\_iff\_avoids\_obstruction}, \texttt{generic\_iff\_meets\_complete\_family}, \texttt{registered\_pairMismatch\_excludes\_duplicate}, \texttt{forcing\_repair\_both} & --- \\
Prop.~\ref{prop:nolens} & \texttt{CohenBoundary}: \texttt{singleton\_visibility\_forces\_injective}, \texttt{singleton\_visibility\_forces\_all} & --- \\
Thm.~\ref{thm:strict}(1)--(2), finite cores & \texttt{CohenBoundary}: \texttt{adjoining\_strict\_iff\_novel}; \texttt{CohenRecognition}: \texttt{generic\_implies\_novel}, \texttt{generic\_implies\_strict\_extension} & --- \\
Prop.~\ref{prop:axioms} & \path{formal/AxiomAudit.lean} (all declarations) & --- \\
\bottomrule
\end{tabularx}
\caption{Results, Lean declarations and diagnostics. Lean statements about Cohen conditions concern
conditions on $\mathbb N$ and explicitly supplied families; their connection to a countable
transitive model is argued on paper.}
\label{tab:artifact-map}
\end{table}
 \section{Imported results}\label{app:extended-material}

Table~\ref{tab:imports-ledger} lists the classical results the paper uses without proof, where
they are used, and what they are not used for.

\begin{table}[htbp]
\centering
\footnotesize
\begin{tabularx}{\textwidth}{@{}>{\raggedright\arraybackslash}p{0.27\textwidth}>{\raggedright\arraybackslash}p{0.30\textwidth}>{\raggedright\arraybackslash}X@{}}
\toprule
Import & Used for & Not used for \\
\midrule
Cumulative hierarchy by transfinite recursion; definability of rank in $\mathsf{ZF}$
\citep{Zermelo1930,Jech2003} & Reading the $U$ of Theorem~\ref{thm:shadow} as the shadow of its own staged hierarchy & Proving Theorem~\ref{thm:shadow}, whose proof is direct \\
Natural models $V_\kappa$ \citep{MontagueVaught1959} & Explaining why truncations need rank-bounded images & Any existence claim for inaccessible cardinals or set models \\
$V_{\omega+\omega}$ and $H_{\omega_1}$ in ambient $\mathsf{ZFC}$ \citep{Jech2003,GitmanHamkinsJohnstone2016} & Non-redundancy of \textbf{C-REPL} and \textbf{C-POW} & A consistency proof for $\mathsf{ZFC}$ \\
Choice equivalences and symmetric models of $\mathsf{ZF}$ \citep{Jech1973,Cohen1966} & Pricing \textbf{C-SEL}; non-redundancy of \textbf{C-SEL} & A free infinitary selector \\
Generic extensions $M[G]$ \citep{Cohen1966,Kunen2011} (\textbf{J2}) & Naming the model containing $x_G$; the $V=L$ example & The novelty, recognition and strictness results, which use \textbf{J1} only \\
Relative consistency of $\mathsf{CH}$ and $\neg\mathsf{CH}$; ccc preservation \citep{Godel1938,Godel1940,Cohen1963,Kunen2011} & Theorem~\ref{thm:ch} & New independence theorems \\
Multiverse and universe views \citep{Hamkins2012,Godel1947} & Discussion & A decision between them \\
\bottomrule
\end{tabularx}
\caption{External imports and their boundaries.}
\label{tab:imports-ledger}
\end{table}
 \section{Version history}\label{app:versions}

\paragraph{v2 (3 October 2026).} The forcing section states the relation between genericity and
novelty as a one-way implication (Theorem~\ref{thm:novelty}), adds an explicit real that is new and
splits every infinite coded set without being generic (Proposition~\ref{prop:converse}), and adds
the exact characterization of genericity by coded closed nowhere dense sets
(Theorem~\ref{thm:recognition}); version~1 had presented genericity as characterized by
non-definability. The ground model's predicates are treated as a Boolean algebra rather than a
partition lens (Proposition~\ref{prop:nolens}, Theorem~\ref{thm:strict}), and the distinction
between category and measure is made explicit (Proposition~\ref{prop:category}). The shadow theorem
states its structural premise (S) explicitly; the non-redundancy witnesses now include
$H_{\omega_1}$ for Power set; countermodel \textbf{CM-1} uses a node-counting lens; and the
continuum-hypothesis statements separate relative consistency from the host-model picture. Smaller
corrections concern the countability of the carrier and of $\powset(\HF)$, the indexing of the
$\HF$ tower, sections of the membership lens, and the hypotheses of G\"odel's theorems. The Lean
development grew from five to eight modules, with an automatic axiom audit over all declarations,
and the diagnostics gained exact-count and statistical gates. The exposition was rewritten
throughout, and figures were added.

\paragraph{v1 (17 June 2026).} First public version.
 
\bibliographystyle{plainnat}
\bibliography{references}

\end{document}
